\section{AMC 12 12B 2018}
        \begin{enumerate}
        	\item (\textit{AMC 12 12B 2018}) Kate bakes a $20$-inch by $18$-inch pan of cornbread. The cornbread is cut into pieces that measure $2$ inches by $2$ inches. How many pieces of cornbread does the pan contain?


 $\textbf{(A) } 90 \qquad \textbf{(B) } 100 \qquad \textbf{(C) } 180 \qquad \textbf{(D) } 200 \qquad \textbf{(E) } 360$

 \textbf{Answer: }A



\textbf{Solution 1}

The area of the pan is $20\cdot18=360$. Since the area of each piece is $2\cdot2=4$, there are $\frac{360}{4} = \boxed{\textbf{(A) } 90}$ pieces.

 



\textbf{Solution 2}

By dividing each of the dimensions by $2$, we get a $10\times9$ grid that makes $\boxed{\textbf{(A) } 90}$ pieces.

 


	\item (\textit{AMC 12 12B 2018}) Sam drove $96$ miles in $90$ minutes. His average speed during the first $30$ minutes was $60$ mph (miles per hour), and his average speed during the second $30$ minutes was $65$ mph. What was his average speed, in mph, during the last $30$ minutes?


 $\textbf{(A) } 64 \qquad \textbf{(B) } 65 \qquad \textbf{(C) } 66 \qquad \textbf{(D) } 67 \qquad \textbf{(E) } 68$

 \textbf{Answer: }D



\textbf{Solution 1}

Suppose that Sam's average speed during the last $30$ minutes was $x$ mph.

 Recall that a half hour is equal to $30$ minutes. Therefore, Sam drove $60\cdot0.5=30$ miles during the first half hour, $65\cdot0.5=32.5$ miles during the second half hour, and $x\cdot0.5$ miles during the last half hour. We have 
 \begin{align*} 30+32.5+x\cdot0.5&=96 \\ x\cdot0.5&=33.5 \\ x&=\boxed{\textbf{(D) } 67}. \end{align*}
~Haha0201 ~MRENTHUSIASM

 



\textbf{Solution 2}

Suppose that Sam's average speed during the last $30$ minutes was $x$ mph.

 Note that Sam's average speed during the entire trip was $\frac{96}{3/2}=64$ mph. Since Sam drove at $60$ mph, $65$ mph, and $x$ mph for the same duration ($30$ minutes each), his average speed during the entire trip was the average of $60$ mph, $65$ mph, and $x$ mph. We have
 \begin{align*} \frac{60+65+x}{3}&=64 \\ 60+65+x&=192 \\ x&=\boxed{\textbf{(D) } 67}. \end{align*}
~coolmath\_2018 ~MRENTHUSIASM

 


	\item (\textit{AMC 12 12B 2018}) A line with slope $2$ intersects a line with slope $6$ at the point $(40,30)$. What is the distance between the $x$-intercepts of these two lines? 


 $\textbf{(A) } 5 \qquad \textbf{(B) } 10 \qquad \textbf{(C) } 20 \qquad \textbf{(D) } 25 \qquad \textbf{(E) } 50$

 \textbf{Answer: }B



\textbf{Solution 1}

Using point slope form, we get the equations $y-30 = 6(x-40)$ and $y-30 = 2(x-40)$. Simplifying, we get $6x-y=210$ and $2x-y=50$. Letting $y=0$ in both equations and solving for $x$ gives the $x$-intercepts: $x=35$ and $x=25$, respectively. Thus the distance between them is $35-25=\boxed{\textbf{(B) } 10}$.

 



\textbf{Solution 2}

In order for the line with slope $2$ to travel "up" $30$ units (from $y=0$), it must have traveled $30/2=15$ units to the right. Thus, the $x$-intercept is at $x=40-15=25$. As for the line with slope $6$, in order for it to travel "up" $30$ units it must have traveled $30/6=5$ units to the right. Thus its $x$-intercept is at $x=40-5=35$. Then the distance between them is $35-25=\boxed{\textbf{(B) } 10}$.

 


	\item (\textit{AMC 12 12B 2018}) A circle has a chord of length $10$, and the distance from the center of the circle to the chord is $5$. What is the area of the circle?


 $\textbf{(A) }25\pi \qquad \textbf{(B) }50\pi \qquad \textbf{(C) }75\pi \qquad \textbf{(D) }100\pi \qquad \textbf{(E) }125\pi \qquad$

 \textbf{Answer: }B



\textbf{Solution}

Let $O$ be the center of the circle, $\overline{AB}$ be the chord, and $M$ be the midpoint of $\overline{AB},$ as shown below.
                    \begin{figure}[H]
                        \centering
                        \includegraphics[width=0.5\textwidth]{Solution/AMC_12_12B_2018_4_Solution_Figure_1.png}
                    \end{figure}
                    Note that $\overline{OM}\perp\overline{AB}.$ Since $OM=AM=BM=5,$ we conclude that $\triangle OMA$ and $\triangle OMB$ are congruent isosceles right triangles. It follows that $r=5\sqrt2,$ so the area of $\odot O$ is $\pi r^2=\boxed{\textbf{(B) }50\pi}$.

 


	\item (\textit{AMC 12 12B 2018}) How many subsets of $\{2,3,4,5,6,7,8,9\}$ contain at least one prime number?


 $\textbf{(A) } 128 \qquad \textbf{(B) } 192 \qquad \textbf{(C) } 224 \qquad \textbf{(D) } 240 \qquad \textbf{(E) } 256$

 \textbf{Answer: }D



\textbf{Solution 1}

Since an element of a subset is either in or out, the total number of subsets of the $8$-element set is $2^8 = 256$. However, since we are only concerned about the subsets with at least $1$ prime in it, we can use complementary counting to count the subsets without a prime and subtract that from the total. Because there are $4$ non-primes, there are $2^8 -2^4 = \boxed{\textbf{(D) } 240}$ subsets with at least $1$ prime.

 



\textbf{Solution 2}

We can construct our subset by choosing which primes are included and which composites are included. There are $2^4-1$ ways to select the primes (total subsets minus the empty set) and $2^4$ ways to select the composites. Thus, there are $15\cdot16=\boxed{\textbf{(D) } 240}$ ways to choose a subset of the eight numbers.

 



\textbf{Solution 3}

We complement count. We know that we have $2^8$ subsets in all, including the empty subset. We subtract $\dbinom{4}{0}+\dbinom{4}{1}+\dbinom{4}{2}+\dbinom{4}{3}+\dbinom{4}{4} = 2^4 = 16$ We have $2^8-16 = 256-16 = \boxed{\textbf{(D) }240}$ ways to choose a subset of the eight numbers that include a prime number. ~hh99754539

 


	\item (\textit{AMC 12 12B 2018}) Suppose $S$ cans of soda can be purchased from a vending machine for $Q$ quarters. Which of the following expressions describes the number of cans of soda that can be purchased for $D$ dollars, where $1$ dollar is worth $4$ quarters?


 $\textbf{(A) } \frac{4DQ}{S} \qquad \textbf{(B) } \frac{4DS}{Q} \qquad \textbf{(C) } \frac{4Q}{DS} \qquad \textbf{(D) } \frac{DQ}{4S} \qquad \textbf{(E) } \frac{DS}{4Q}$

 \textbf{Answer: }B



\textbf{Solution 1}

Each can of soda costs $\frac QS$ quarters, or $\frac{Q}{4S}$ dollars. Therefore, $D$ dollars can purchase $\frac{D}{\left(\frac{Q}{4S}\right)}=\boxed{\textbf{(B) } \frac{4DS}{Q}}$ cans of soda.

 



\textbf{Solution 2}

Note that $S$ is in the unit of $\text{can}.$ On the other hand, $Q$ and $D$ are both in the unit of $\text{cost}.$

 The units of $\textbf{(A)},\textbf{(B)},\textbf{(C)},\textbf{(D)},$ and $\textbf{(E)}$ are $\frac{\text{cost}^2}{\text{can}},\text{can},\frac{1}{\text{can}},\frac{\text{cost}^2}{\text{can}},$ and $\text{can},$ respectively. Since the answer is in the unit of $\text{can},$ we eliminate $\textbf{(A)},\textbf{(C)},$ and $\textbf{(D)}.$ Moreover, it is clear that $D$ dollars can purchase more than $S=\frac{DS}{4Q}$ cans of soda, so we eliminate $\textbf{(E)}.$

 Finally, the answer is $\boxed{\textbf{(B) } \frac{4DS}{Q}}.$



	\item (\textit{AMC 12 12B 2018}) What is the value of 
 \[\log_37\cdot\log_59\cdot\log_711\cdot\log_913\cdots\log_{21}25\cdot\log_{23}27?\]
$\textbf{(A) } 3 \qquad \textbf{(B) } 3\log_{7}23 \qquad \textbf{(C) } 6 \qquad \textbf{(D) } 9 \qquad \textbf{(E) } 10$

 \textbf{Answer: }C



\textbf{Solution 1}

From the Change of Base Formula, we have 
 \[\frac{\prod_{i=3}^{13} \log (2i+1)}{\prod_{i=1}^{11}\log (2i+1)} = \frac{\log 25 \cdot \log 27}{\log 3 \cdot \log 5} = \frac{(2\log 5)\cdot(3\log 3)}{\log 3 \cdot \log 5} = \boxed{\textbf{(C) } 6}.\]




\textbf{Solution 2}

Using the chain rule of logarithms $\log _{a} b \cdot \log _{b} c = \log _{a} c,$ we get
 \begin{align*} \log_37\cdot\log_59\cdot\log_711\cdot\log_913\cdots\log_{21}25\cdot\log_{23}27 &= (\log _{3} 7 \cdot \log _{7} 11 \cdots \log _{23} 27) \cdot (\log _{5} 9 \cdot \log _{9} 13 \cdots \log _{21} 25) \\ &= \log _{3} 27 \cdot \log _{5} 25 \\ &= 3 \cdot 2 \\ &= \boxed{\textbf{(C) } 6}. \end{align*}




	\item (\textit{AMC 12 12B 2018}) Line segment $\overline{AB}$ is a diameter of a circle with $AB = 24$. Point $C$, not equal to $A$ or $B$, lies on the circle. As point $C$ moves around the circle, the centroid (center of mass) of $\triangle ABC$ traces out a closed curve missing two points. To the nearest positive integer, what is the area of the region bounded by this curve?


 $\textbf{(A) } 25 \qquad \textbf{(B) } 38  \qquad \textbf{(C) } 50  \qquad \textbf{(D) } 63 \qquad \textbf{(E) } 75$

 \textbf{Answer: }C



\textbf{Solution 1}

By the Inscribed Angle Theorem, $\triangle ABC$ is a right triangle with $\angle C=90^{\circ}.$ So, its circumcenter is the midpoint of $\overline{AB},$ and its median from $C$ is half as long as $\overline{AB}.$ For each $\triangle ABC,$ let $O$ and $G$ be its circumcenter and centroid, respectively. It follows that $OA=OB=OC=12.$ In any triangle, since the centroid divides each median into parts in the ratio $2:1,$ with the centroid being twice as close to the midpoint of a side as it is to the opposite vertex, we have $OG=\frac13 OC=4.$

 As shown below, $\triangle ABC_1$ and $\triangle ABC_2$ are two shapes of $\triangle ABC$ with centroids $G_1$ and $G_2,$ respectively:
                    \begin{figure}[H]
                        \centering
                        \includegraphics[width=0.5\textwidth]{Solution/AMC_12_12B_2018_8_Solution_Figure_1.png}
                    \end{figure}
                    Therefore, point $G$ traces out a circle (missing two points) with the center $O$ and the radius $\overline{OG},$ as indicated in red. To the nearest positive integer, the area of the region bounded by the red curve is $\pi\cdot OG^2=16\pi\approx\boxed{\textbf{(C) } 50}.$




\textbf{Solution 2}

We assign coordinates. Let $A = (-12,0)$, $B = (12,0)$, and $C = (x,y)$ lie on the circle $x^2 +y^2 = 12^2$. Then, the centroid of $\triangle ABC$ is $G = \left(\frac{-12 + 12 + x}{3}, \frac{0 + 0 + y}{3}\right) = \left(\frac x3,\frac y3\right)$. Thus, $G$ traces out a circle with a radius $\frac13$ of the radius of the circle that point $C$ travels on. Thus, $G$ traces out a circle of radius $\frac{12}{3} = 4$, which has area $16\pi\approx \boxed{\textbf{(C) } 50}$.

 



\textbf{Solution 3}

First we can draw a few conclusions from the given information. Firstly we can see clearly that the distance from the centroid to the center of the circle will remain the same no matter $C$ is on the circle. Also we can see that because the two legs of every triangles will always originate on the diameter, using inscribed angle rules, we know that $\angle C = \frac{180^\circ}{2} = 90^\circ$. Now we know that all triangles $ABC$ will be a right triangle. We also know that the closed curve will simply be a circle with radius equal to the centroid of each triangle. We can now pick any arbitrary triangle, calculate its centroid, and the plug it in to the area formula. Using a $45^\circ$-$45^\circ$-$90^\circ$ triangle in conjunction with the properties of a centroid, we can quickly see that the length of the centroid is $4$ now we can plug it in to the area formula where we get $16\pi\approx\boxed{\textbf{(C) } 50}$.

 


 



\textbf{Solution 4}

Let the center of the circle be $O$. Since $AB$ is a diameter of length $24$, $OC$ is a median to $AB$ and a radius of the circle, so $OC = 12$. 

 Let the centroid of $\triangle ABC$ be $G$. By the Centroid Theorem, $G$ lies on the segment $OC$ such that the ratio of $OG$ to $GC$ is $1:2$. 

 Since $OG = \frac{1}{3} OC$, we have:
 \[OG = \frac{1}{3}(12) = 4\]

 As $C$ moves along the boundary of the circle, the centroid $G$ carves out a smaller, concentric circle of radius $4$. 

 The area of this region is:
 \[\text{Area} = \pi \cdot 4^2 = 16\pi \approx 50.27\]

 Rounding to the nearest whole number, the answer is $\boxed{50}$.

 



\textbf{Solution 5}

By barycentric coordinates, $G$ the centroid satisfies $G = \frac{a+b+c}{3}$ where $a,b,c$ are the position vectors of $A,B,C$ respectively. Then, by definition, let $a \langle 0,0 \rangle$ and $b \langle 24, 0 \rangle$. Then, $G = \langle 8,0 \rangle + \frac{c}{3}$. As $C$ lies on the circle, the vector value of $G$ is solely dependent on $\frac{c}{3}$. That is, as $C$ can lie on the circle, $\frac{c}{3}$ oscillates between $\langle 8,0 \rangle$ and $\langle 16,0 \rangle$, on the circle. Anywhere $C$ traverses, that value of $G$ is correspondent, hence if $C$ moves along the circle, then $G$ does as well along that interval. The requested path is then a circle concentric to circle $ABC$ whose radius is 4. That area is then $16 \cdot \pi \approx 16 \cdot 3 = 48$, so the nearest answer choice is $\boxed{50}$.

 


	\item (\textit{AMC 12 12B 2018}) What is

 \[\sum^{100}_{i=1} \sum^{100}_{j=1} (i+j) ?\]
$\textbf{(A) }100{,}100 \qquad \textbf{(B) }500{,}500\qquad \textbf{(C) }505{,}000 \qquad \textbf{(D) }1{,}001{,}000 \qquad \textbf{(E) }1{,}010{,}000 \qquad$

 \textbf{Answer: }E



\textbf{Solution 1}

Recall that the sum of the first $100$ positive integers is $\sum^{100}_{k=1} k = \frac{101\cdot100}{2}=5050.$ It follows that
 \begin{align*} \sum^{100}_{i=1} \sum^{100}_{j=1} (i+j) &= \sum^{100}_{i=1} \sum^{100}_{j=1}i + \sum^{100}_{i=1} \sum^{100}_{j=1}j \\ &= \sum^{100}_{i=1} (100i) + 100 \sum^{100}_{j=1}j \\ &= 100 \sum^{100}_{i=1}i + 100 \sum^{100}_{j=1}j \\ &= 100\cdot5050 + 100\cdot5050 \\ &= \boxed{\textbf{(E) }1{,}010{,}000}. \end{align*}
~MRENTHUSIASM

 



\textbf{Solution 2}

Recall that the sum of the first $100$ positive integers is $\sum^{100}_{k=1} k = \frac{101\cdot100}{2}=5050.$ It follows that
 \begin{align*} \sum^{100}_{i=1} \sum^{100}_{j=1} (i+j) &= \sum^{100}_{i=1} \left(\sum^{100}_{j=1} i + \sum^{100}_{j=1} j\right) \\ &= \sum^{100}_{i=1} (100i+5050) \\ &= 100\sum^{100}_{i=1} i + \sum^{100}_{i=1} 5050 \\ &= 100\cdot5050+5050\cdot100 \\ &= \boxed{\textbf{(E) }1{,}010{,}000}. \end{align*}
~Vfire ~MRENTHUSIASM

 



\textbf{Solution 3}

Recall that the sum of the first $100$ positive integers is $\sum^{100}_{k=1} k = \frac{101\cdot100}{2}=5050.$ Since the nested summation is symmetric with respect to $i$ and $j,$ it follows that
 \begin{align*} \sum^{100}_{i=1} \sum^{100}_{j=1} (i+j) &= \sum^{100}_{i=1} \sum^{100}_{i=1} (2i) \\ &= 2\sum^{100}_{i=1} \sum^{100}_{i=1} i \\ &= 2\sum^{100}_{i=1} 5050 \\ &= 2\cdot(5050\cdot100) \\ &= \boxed{\textbf{(E) }1{,}010{,}000}. \end{align*}
~Vfire ~MRENTHUSIASM

 



\textbf{Solution 4}

The sum contains $100\cdot100=10000$ terms, and the average value of both $i$ and $j$ is $\frac{101}{2}.$ Therefore, the sum becomes 
 \[10000\cdot\left(\frac{101}{2}+\frac{101}{2}\right)=\boxed{\textbf{(E) }1{,}010{,}000}.\]~Rejas ~MRENTHUSIASM

 



\textbf{Solution 5}

We start by writing out the first few terms:
 \[\begin{array}{ccccccccc} (1+1) &+ &(1+2) &+ &(1+3) &+ &\cdots &+ &(1+100) \\ (2+1) &+ &(2+2) &+ &(2+3) &+ &\cdots &+ &(2+100) \\ (3+1) &+ &(3+2) &+ &(3+3) &+ &\cdots &+ &(3+100) \\ [-1ex] &&&&\vdots&&&& \\ (100+1) &+ &(100+2) &+ &(100+3) &+ &\cdots &+ &(100+100). \end{array}\]From the first terms in the parentheses, the sum $1+2+3+\cdots+100$ occurs $100$ times vertically.

 From the second terms in the parentheses, the sum $1+2+3+\cdots+100$ occurs $100$ times horizontally.

 Recall that the sum of the first $100$ positive integers is $1+2+3+\cdots+100=\frac{101\cdot100}{2}=5050.$ Therefore, the answer is 
 \[2\cdot\left(5050\cdot100\right)=\boxed{\textbf{(E) }1{,}010{,}000}.\]~RandomPieKevin ‎~MRENTHUSIASM

 



\textbf{Solution 6}

When we expand the nested summation as shown in Solution 5, note that:

 
            \begin{enumerate} 
        \item The term $2$ occurs $1$ time. The term $3$ occurs $2$ times. 

 The term $4$ occurs $3$ times. 

 $\cdots$ 

 The term $101$ occurs $100$ times. 

 More generally, the term $k+1$ occurs $k$ times for $k\in\{1,2,3,\ldots,100\}.$
 
        \item The term $102$ occurs $99$ times. The term $103$ occurs $98$ times. 

 The term $104$ occurs $97$ times. 

 $\cdots$ 

 The term $200$ occurs $1$ time. 

 More generally, the term $k+101$ occurs $100-k$ times for $k\in\{1,2,3,\ldots,99\}.$

            \end{enumerate}
            Together, the nested summation becomes 
 \begin{align*} \sum^{100}_{k=1}\left[(k+1)k\right] + \sum^{99}_{k=1}\left[(k+101)(100-k)\right] &= \sum^{100}_{k=1}\left[k^2+k\right] + \sum^{99}_{k=1}\left[-k^2-k+10100\right] \\ &= \sum^{100}_{k=1}k^2 + \sum^{100}_{k=1}k - \sum^{99}_{k=1}k^2 - \sum^{99}_{k=1}k + \sum^{99}_{k=1}10100 \\ &= \left(\sum^{100}_{k=1}k^2 - \sum^{99}_{k=1}k^2\right) + \left(\sum^{100}_{k=1}k - \sum^{99}_{k=1}k\right) + \sum^{99}_{k=1}10100 \\ &= 100^2+100+10100\cdot99 \\ &= \boxed{\textbf{(E) }1{,}010{,}000}. \end{align*}
~MRENTHUSIASM

 


	\item (\textit{AMC 12 12B 2018}) A list of $2018$ positive integers has a unique mode, which occurs exactly $10$ times. What is the least number of distinct values that can occur in the list?


 $\textbf{(A) }202 \qquad \textbf{(B) }223 \qquad \textbf{(C) }224 \qquad \textbf{(D) }225 \qquad \textbf{(E) }234 \qquad$

 \textbf{Answer: }D



\textbf{Solution 1 (Statistics)}

To minimize the number of distinct values, we want to maximize the number of times a number appears. So, we could have $223$ numbers appear $9$ times, $1$ number appear once, and the mode appear $10$ times, giving us a total of $223 + 1 + 1 = \boxed{\textbf{(D)}\ 225}.$




\textbf{Solution 2 (Algebra)}

As in Solution 1, we want to maximize the number of time each number appears to do so. We can set up an equation $10 + 9( x - 1 )\geq2018,$ where $x$ is the number of values. Notice how we can then rearrange the equation into $1 + 9 ( 1 )+9 ( x - 1 )\geq2018,$ which becomes $9 x\geq2017,$ or $x\geq224\frac19.$ We cannot have a fraction of a value so we must round up to $\boxed{\textbf{(D)}\ 225}.$




\textbf{Solution 3}

There is a unique number \( n \) with a mode of 10. This means we have 2008 terms left. We see that \( \lfloor 2008/9 \rfloor = 223 \), with remainder 1. Our answer is simply \( 223 + 1 + 1 =\) $\boxed{\textbf{(D)}\ 225}.$



	\item (\textit{AMC 12 12B 2018}) A closed box with a square base is to be wrapped with a square sheet of wrapping paper. The box is centered on the wrapping paper with the vertices of the base lying on the midlines of the square sheet of paper, as shown in the figure on the left. The four corners of the wrapping paper are to be folded up over the sides and brought together to meet at the center of the top of the box, point $A$ in the figure on the right. The box has base length $w$ and height $h$. What is the area of the sheet of wrapping paper?


 
                    \begin{figure}[H]
                        \centering
                        \includegraphics[width=0.5\textwidth]{Problem/AMC_12_12B_2018_11_Problem_Figure_1.png}
                    \end{figure}
                    
$\textbf{(A) } 2(w+h)^2 \qquad \textbf{(B) } \frac{(w+h)^2}2 \qquad \textbf{(C) } 2w^2+4wh \qquad \textbf{(D) } 2w^2 \qquad \textbf{(E) } w^2h$

 \textbf{Answer: }A



\textbf{Solution 0, super simple and quick}

You could see that half the diagonal of the square wrapping paper is 
 \[\frac{w}{2} + h + \frac{w}{2} = w + h\]. Thus the full diagonal of the square wrapping paper would be 
 \[d = 2(w + h)\]. So the area would be 
 \[A = \frac{d^2}{2} = \frac{[2(w + h)]^2}{2} = \frac{4(w + h)^2}{2} = 2(w + h)^2\]. Thus the answer is 
 \[\boxed{\text{(A) } 2(w+h)^2}\]




\textbf{Solution 1}

Consider one-quarter of the image (the wrapping paper is divided up into $4$ congruent squares). The length of each dotted line is $h$. The area of the rectangle that is $w$ by $h$ is $wh$. The combined figure of the two triangles with base $h$ is a square with $h$ as its diagonal. Using the Pythagorean Theorem, each side of this square is $\frac{h}{\sqrt2}$. Thus, the area is the side length squared which is $\frac{h^2}{2}$. Similarly, the combined figure of the two triangles with base $w$ is a square with area $\frac{w^2}{2}$. Adding all of these together, we get $\frac{w^2}{2} + \frac{h^2}{2} + wh$. Since we have four of these areas in the entire wrapping paper, we multiply this by $4$, getting $4\left(\frac{w^2}{2} + \frac{h^2}{2} + wh\right) = 2\left(w^2 + h^2 + 2wh\right) = \boxed{\textbf{(A) } 2(w+h)^2}$.

 The diagram for this solution is shown below:
                    \begin{figure}[H]
                        \centering
                        \includegraphics[width=0.5\textwidth]{Solution/AMC_12_12B_2018_11_Solution_Figure_1.png}
                    \end{figure}
                    ~Hydroquantum (Solution)

 



\textbf{Solution 2}

The sheet of paper is made out of $4$ squares, as shown in the diagram in Solution 1. Each square has a side length of $\frac{w}{\sqrt{2}} + \frac{h}{\sqrt{2}}$, which we get from the Pythagorean Theorem (a $45^\circ\text{-}45^\circ\text{-}90^\circ$ triangle's legs is the hypotenuse divided by $\sqrt2$). Thus, to find the area of the entire paper, we square our side length and multiply by $4$. So, the answer is $4\cdot\left(\frac{w}{\sqrt{2}} + \frac{h}{\sqrt{2}}\right)^2 = \boxed{\textbf{(A) } 2(w+h)^2}$.

 



\textbf{Solution 3}

The sheet of paper is made out of the surface area of the box plus the sum of the four yellow triangles, as shown below. 

 
                    \begin{figure}[H]
                        \centering
                        \includegraphics[width=0.5\textwidth]{Solution/AMC_12_12B_2018_11_Solution_Figure_2.png}
                    \end{figure}
                    

 The surface area is $2w^2 + 2wh + 2wh$ which equals $2w^2 + 4wh$.The four triangles each have a height and a base of $h$, so they each have an area of $\frac{h^2}{2}$. There are four of them, so multiplied by four is $2h^2$. Together, paper's area is $2w^2 + 4wh + 2h^2$. This can be factored and written as $\boxed{\textbf{(A) } 2(w+h)^2} \qquad$.

 


	\item (\textit{AMC 12 12B 2018}) Side $\overline{AB}$ of $\triangle ABC$ has length $10$. The bisector of angle $A$ meets $\overline{BC}$ at $D$, and $CD = 3$. The set of all possible values of $AC$ is an open interval $(m,n)$. What is $m+n$?


 $\textbf{(A) }16 \qquad \textbf{(B) }17 \qquad \textbf{(C) }18 \qquad \textbf{(D) }19 \qquad \textbf{(E) }20 \qquad$

 \textbf{Answer: }C



\textbf{Solution}

Let $AC=x.$ By Angle Bisector Theorem, we have $\frac{AB}{AC}=\frac{BD}{CD},$ from which $BD=CD\cdot\frac{AB}{AC}=\frac{30}{x}.$

 Recall that $x>0.$ We apply the Triangle Inequality to $\triangle ABC:$

            \begin{enumerate} 
        \item $AC+BC>AB \iff x+\left(\frac{30}{x}+3\right)>10$ We simplify and complete the square to get $\left(x-\frac72\right)^2+\frac{71}{4}>0,$ from which $x>0.$


        \item $AB+BC>AC \iff 10+\left(\frac{30}{x}+3\right)>x$ We simplify and factor to get $(x+2)(x-15)<0,$ from which $0<x<15.$


        \item $AB+AC>BC \iff 10+x>\frac{30}{x}+3$ We simplify and factor to get $(x+10)(x-3)>0,$ from which $x>3.$


            \end{enumerate}
            Taking the intersection of the solutions gives 
 \[(m,n)=(0,\infty)\cap(0,15)\cap(3,\infty)=(3,15),\] so the answer is $m+n=\boxed{\textbf{(C) }18}.$



	\item (\textit{AMC 12 12B 2018}) Square $ABCD$ has side length $30$. Point $P$ lies inside the square so that $AP = 12$ and $BP = 26$. The centroids of $\triangle{ABP}$, $\triangle{BCP}$, $\triangle{CDP}$, and $\triangle{DAP}$ are the vertices of a convex quadrilateral. What is the area of that quadrilateral? 


 
                    \begin{figure}[H]
                        \centering
                        \includegraphics[width=0.5\textwidth]{Problem/AMC_12_12B_2018_13_Problem_Figure_1.png}
                    \end{figure}
                    
$\textbf{(A) }100\sqrt{2}\qquad\textbf{(B) }100\sqrt{3}\qquad\textbf{(C) }200\qquad\textbf{(D) }200\sqrt{2}\qquad\textbf{(E) }200\sqrt{3}$

 \textbf{Answer: }C



\textbf{Solution 1 (Similar Triangles)}

As shown below, let $M_1,M_2,M_3,M_4$ be the midpoints of $\overline{AB},\overline{BC},\overline{CD},\overline{DA},$ respectively, and $G_1,G_2,G_3,G_4$ be the centroids of $\triangle{ABP},\triangle{BCP},\triangle{CDP},\triangle{DAP},$ respectively.
                    \begin{figure}[H]
                        \centering
                        \includegraphics[width=0.5\textwidth]{Solution/AMC_12_12B_2018_13_Solution_Figure_1.png}
                    \end{figure}
                    By SAS, we conclude that $\triangle G_1G_2P\sim\triangle M_1M_2P, \triangle G_2G_3P\sim\triangle M_2M_3P, \triangle G_3G_4P\sim\triangle M_3M_4P,$ and $\triangle G_4G_1P\sim\triangle M_4M_1P.$ By the properties of centroids, the ratio of similitude for each pair of triangles is $\frac{2}{3}.$

 Note that quadrilateral $M_1M_2M_3M_4$ is a square of side-length $15\sqrt2.$ It follows that:

 
            \begin{enumerate} 
        \item Since $\overline{G_1G_2}\parallel\overline{M_1M_2},\overline{G_2G_3}\parallel\overline{M_2M_3},\overline{G_3G_4}\parallel\overline{M_3M_4},$ and $\overline{G_4G_1}\parallel\overline{M_4M_1}$ by the Converse of the Corresponding Angles Postulate, we have $\angle G_1G_2G_3=\angle G_2G_3G_4=\angle G_3G_4G_1=\angle G_4G_1G_2=90^\circ.$

        \item Since $G_1G_2=\frac23M_1M_2, G_2G_3=\frac23M_2M_3, G_3G_4=\frac23M_3M_4,$ and $G_4G_1=\frac23M_4M_1$ by the ratio of similitude, we have $G_1G_2=G_2G_3=G_3G_4=G_4G_1=10\sqrt2.$

            \end{enumerate}
            Together, quadrilateral $G_1G_2G_3G_4$ is a square of side-length $10\sqrt2,$ so its area is $\left(10\sqrt2\right)^2=\boxed{\textbf{(C) }200}.$
<u>\textbf{Remark}</u>

 This solution shows that, if point $P$ is within square $ABCD,$ then the shape and the area of quadrilateral $G_1G_2G_3G_4$ are independent of the location of $P.$ Let the brackets denote areas. More generally, $G_1G_2G_3G_4$ is always a square of area 
 \[[G_1G_2G_3G_4]=\left(\frac23\right)^2[M_1M_2M_3M_4]=\frac49[M_1M_2M_3M_4]=\frac29[ABCD].\] On the other hand, the location of $G_1G_2G_3G_4$ is dependent on the location of $P.$




\textbf{Solution 2 (Similar Triangles)}

This solution refers to the diagram in Solution 1.

 By SAS, we conclude that $\triangle G_1G_3P\sim\triangle M_1M_3P$ and $\triangle G_2G_4P\sim\triangle M_2M_4P.$ By the properties of centroids, the ratio of similitude for each pair of triangles is $\frac23.$

 Note that quadrilateral $M_1M_2M_3M_4$ is a square of diagonal-length $30,$ so $\overline{M_1M_3}\perp\overline{M_2M_4}.$ Since $\overline{G_1G_3}\parallel\overline{M_1M_3}$ and $\overline{G_2G_4}\parallel\overline{M_2M_4}$ by the Converse of the Corresponding Angles Postulate, we have $\overline{G_1G_3}\perp\overline{G_2G_4}.$

 Therefore, the area of quadrilateral $G_1G_2G_3G_4$ is 
 \[\frac12\cdot G_1G_3\cdot G_2G_4 = \frac12\cdot\left(\frac23\cdot M_1M_3\right)\cdot\left(\frac23\cdot M_2M_4\right)=\boxed{\textbf{(C) }200}.\]~Funnybunny5246 ~MRENTHUSIASM

 



\textbf{Solution 3 (Coordinate Geometry)}

This solution refers to the diagram in Solution 1.

 We place the diagram in the coordinate plane: Let $A=(0,30),B=(0,0),C=(30,0),D=(30,30),$ and $P=(3x,3y).$

 Recall that for any triangle in the coordinate plane, the coordinates of its centroid are the averages of the coordinates of its vertices. It follows that $G_1=(x,y+10),G_2=(x+10,y),G_3=(x+20,y+10),$ and $G_4=(x+10,y+20).$

 Note that $G_1G_3=G_2G_4=20$ and $\overline{G_1G_3}\perp\overline{G_2G_4}.$ Therefore, the area of quadrilateral $G_1G_2G_3G_4$ is 
 \[\frac12\cdot G_1G_3\cdot G_2G_4=\boxed{\textbf{(C) }200}.\]




\textbf{Solution 4 (Homothety)}

Let $X,Y,Z,W$ be the midpoints of sides $AB,BC,CD,DE$, respectively.

 Notice that a homothety centered at P with ratio $\frac{2}{3}$ will send $XYZW$ to $G_{1}G_{2}G_{3}G_{4}$, so $G_{1}G_{2}G_{3}G_{4}$ is a square with area $\left(\frac{2}{3}\right)^2 [XYZW]$, but $[XYZW]=\frac{1}{2}[ABCD]$ so our desired area is 
 \[\frac{4}{9}\cdot\frac{1}{2}\cdot900=\boxed{\textbf{(C) }200}\]

 Desmos graph:\href{https://www.desmos.com/calculator/cqozcf5ef1}{https://www.desmos.com/calculator/cqozcf5ef1}~Ynsg

 


	\item (\textit{AMC 12 12B 2018}) Joey and Chloe and their daughter Zoe all have the same birthday. Joey is $1$ year older than Chloe, and Zoe is exactly $1$ year old today. Today is the first of the $9$ birthdays on which Chloe's age will be an integral multiple of Zoe's age. What will be the sum of the two digits of Joey's age the next time his age is a multiple of Zoe's age?


 $\textbf{(A) }7 \qquad \textbf{(B) }8 \qquad \textbf{(C) }9 \qquad \textbf{(D) }10 \qquad \textbf{(E) }11 \qquad$

 \textbf{Answer: }E



\textbf{Solution 1}

Suppose that Chloe is $c$ years old today, so Joey is $c+1$ years old today. After $n$ years, Chloe and Zoe will be $n+c$ and $n+1$ years old, respectively. We are given that 
 \[\frac{n+c}{n+1}=1+\frac{c-1}{n+1}\] is an integer for $9$ nonnegative integers $n.$ It follows that $c-1$ has $9$ positive divisors. The prime factorization of $c-1$ is either $p^8$ or $p^2q^2.$ Since $c-1<100,$ the only possibility is $c-1=2^2\cdot3^2=36,$ from which $c=37.$ We conclude that Joey is $c+1=38$ years old today.

 Suppose that Joey's age is a multiple of Zoe's age after $k$ years, in which Joey and Zoe will be $k+38$ and $k+1$ years old, respectively. We are given that 
 \[\frac{k+38}{k+1}=1+\frac{37}{k+1}\] is an integer for some positive integer $k.$ It follows that $37$ is divisible by $k+1,$ so the only possibility is $k=36.$ We conclude that Joey will be $k+38=74$ years old then, from which the answer is $7+4=\boxed{\textbf{(E) }11}.$




\textbf{Solution 2}

Let Joey's age be $j$, Chloe's age be $c$, and we know that Zoe's age is $1$. 

 We know that there must be $9$ values $k\in\mathbb{Z}$ such that $c+k=a(1+k)$ where $a$ is an integer.

 Therefore, $c-1+(1+k)=a(1+k)$ and $c-1=(1+k)(a-1)$. Therefore, we know that, as there are $9$ solutions for $k$, there must be $9$ solutions for $c-1$. We know that this must be a perfect square. Testing perfect squares, we see that $c-1=36$, so $c=37$. Therefore, $j=38$. Now, since $j-1=37$, by similar logic, $37=(1+k)(a-1)$, so $k=36$ and Joey will be $38+36=74$ and the sum of the digits is $\boxed{\textbf{(E) }11}$.

 



\textbf{Solution 3}

Let Chloe's age be $c$ so that Zoey is 1. Then, if $n$ years pass, we have that again $c+n \equiv 0 \pmod{n+1}$ from which $c \equiv 1 \pmod{n+1}$ thus $c - 1 \equiv 0 \pmod{n+1}$. Hence, $c-1$ has 9 divisors. Assuming no human lives past a century, it is given that only numbers of the form $p^2q^2$ work. Taking minimality, $p = 2, q = 3$ from which $c-1 = 36 \implies c = 37$. Hence, Joey is 38. Now, if after some $k$ years Joey is a multiple of Zoe, we can just apply the same process as before. In other words, $38+k \equiv 0 \pmod{1+k}$ hence $38 \equiv 1 \pmod{1+k}$, and thus $37 \equiv 0 \pmod{1+k}$. The only integer multiples of 37 is 1 and 37, but they wish for the next integral multiple, which is when Zoe is 38 years and Joey is 74, whose sum is $\boxed{11}$.

 



\textbf{Solution 4}

Here's a different way of stating Solution 2:

 If a number is a multiple of both Chloe's age and Zoe's age, then it is a multiple of their difference. Since the difference between their ages does not change, then that means the difference between their ages has $9$ factors. Therefore, the difference between Chloe and Zoe's age is $36$, so Chloe is $37$, and Joey is $38$. The common factor that will divide both of their ages is $37$, so Joey will be $74$. The answer is $7 + 4 =  \boxed{\textbf{(E) }11}$.

 



\textbf{Solution 5}

Similar approach to above, just explained less concisely and more in terms of the problem (less algebraic).

 Let $C+n$ denote Chloe's age, $J+n$ denote Joey's age, and $Z+n$ denote Zoe's age, where $n$ is the number of years from now. We are told that $C+n$ is a multiple of $Z+n$ exactly nine times. Because $Z+n$ is $1$ at $n=0$ and will increase until greater than $C-Z$, it will hit every natural number less than $C-Z$, including every factor of $C-Z$. For $C+n$ to be an integral multiple of $Z+n$, the difference $C-Z$ must also be a multiple of $Z$, which happens if $Z$ is a factor of $C-Z$. Therefore, $C-Z$ has nine factors. The smallest number that has nine positive factors is $2^23^2=36$ . (We want it to be small so that Joey will not have reached three digits of age before his age is a multiple of Zoe's.) We also know $Z=1$ and $J=C+1$. Thus,
 \begin{align*} C-Z&=36, \\ J-Z&=37. \end{align*}
By our above logic, the next time $J-Z$ is a multiple of $Z+n$ will occur when $Z+n$ is a factor of $J-Z$. Because $37$ is prime, the next time this happens is at $Z+n=37$, when $J+n=74$. The answer is $7+4=\boxed{\textbf{(E) }11}$.

 



\textbf{Solution 6 (Number Theory+Logic)}

Let Joeys current age be $J$, Chloe's current age be $C$, Zoe's current age be $Z$, and $C-Z=C-1=K$. This gives $C=1+K$ and $J=2+K$. After $t$ years, Zoe will be $Z+t=t+1$ years old and Chloe will be $C+t$ years old. Chloe's age is an integral multiple of Zoe's age when $Z+t|C+t$. By the Divisibility of Differences rule (which is a unique case of the linear combination property of divisibility), $t+1|C+t-(t+1)$->$t+1|C-1$->$t+1=K$.

 $K$ must have exactly $9$ positive integer factors, since there are precisely $9$ values of $t$ for which Zoe's age is a factor of Chloe's age. The smallest number with this property is $2^2 \times 3^2=36$. We assume this is the value of $K$, as the next smallest value of $K$ (which is $K=100$) yields the unlikely case that Chloe is currently $101$ years old. 

 Since $K=36$, we have $C=1+36=37$ and $J=2+36=38$. Joey's age is a multiple of Zoe's age when $Z+t|J+t$. Plugging in the values of $Z$ and $J$, we get $t+1|t+38$. By the Divisibility of Differences rule, $t+1|t+38-(t+1)=37$. The smallest value of $t$ is $36$, so Joey will be $38+36=74$ years old. The sum of the digits is $7+4=\boxed{\textbf{(E) } 11}$.

 


	\item (\textit{AMC 12 12B 2018}) How many odd positive $3$-digit integers are divisible by $3$ but do not contain the digit $3$?


 $\textbf{(A) } 96 \qquad \textbf{(B) } 97 \qquad \textbf{(C) } 98 \qquad \textbf{(D) } 102 \qquad \textbf{(E) } 120$

 \textbf{Answer: }A



\textbf{Solution 1}

Let $\underline{ABC}$ be one such odd positive $3$-digit integer with hundreds digit $A,$ tens digit $B,$ and ones digit $C.$ Since $\underline{ABC}\equiv0\pmod3,$ we need $A+B+C\equiv0\pmod3$ by the divisibility rule for $3.$

 As $A\in\{1,2,4,5,6,7,8,9\}$ and $C\in\{1,5,7,9\},$ there are $8$ possibilities for $A$ and $4$ possibilities for $C.$ Note that each ordered pair $(A,C)$ determines the value of $B$ modulo $3,$ so $B$ can be any element in one of the sets $\{0,6,9\},\{1,4,7\},$ or $\{2,5,8\}.$ We conclude that there are always $3$ possibilities for $B.$

 By the Multiplication Principle, the answer is $8\cdot4\cdot3=\boxed{\textbf{(A) } 96}.$




\textbf{Solution 2}

Let $\underline{ABC}$ be one such odd positive $3$-digit integer with hundreds digit $A,$ tens digit $B,$ and ones digit $C.$ Since $\underline{ABC}\equiv0\pmod3,$ we need $A+B+C\equiv0\pmod3$ by the divisibility rule for $3.$

 As $A\in\{1,2,4,5,6,7,8,9\},B\in\{0,1,2,4,5,6,7,8,9\},$ and $C\in\{1,5,7,9\},$ note that:

 
            \begin{enumerate} 
        \item There are $2$ possibilities for $A\equiv0\pmod3,$ namely $A=6,9.$ There are $3$ possibilities for $A\equiv1\pmod3,$ namely $A=1,4,7.$ 

 There are $3$ possibilities for $A\equiv2\pmod3,$ namely $A=2,5,8.$
 
        \item There are $3$ possibilities for $B\equiv0\pmod3,$ namely $B=0,6,9.$ There are $3$ possibilities for $B\equiv1\pmod3,$ namely $B=1,4,7.$ 

 There are $3$ possibilities for $B\equiv2\pmod3,$ namely $B=2,5,8.$
 
        \item There are $1$ possibility for $C\equiv0\pmod3,$ namely $C=9.$ There are $2$ possibilities for $C\equiv1\pmod3,$ namely $C=1,7.$ 

 There are $1$ possibility for $C\equiv2\pmod3,$ namely $C=5.$

            \end{enumerate}
            We apply casework to $A+B+C\equiv0\pmod3:$
 \[\begin{array}{c|c|c||l} & & & \\ [-2.5ex] \boldsymbol{A\operatorname{mod}3} & \boldsymbol{B\operatorname{mod}3} & \boldsymbol{C\operatorname{mod}3} & \multicolumn{1}{c}{\textbf{Count}} \\ [0.5ex] \hline & & &  \\ [-2ex] 0 & 0 & 0 & 2\cdot3\cdot1=6 \\ 0 & 1 & 2 & 2\cdot3\cdot1=6 \\ 0 & 2 & 1 & 2\cdot3\cdot2=12 \\ 1 & 0 & 2 & 3\cdot3\cdot1=9 \\ 1 & 1 & 1 & 3\cdot3\cdot2=18 \\ 1 & 2 & 0 & 3\cdot3\cdot1=9 \\ 2 & 0 & 1 & 3\cdot3\cdot2=18 \\ 2 & 1 & 0 & 3\cdot3\cdot1=9 \\ 2 & 2 & 2 & 3\cdot3\cdot1=9 \end{array}\]Together, the answer is $6+6+12+9+18+9+18+9+9=\boxed{\textbf{(A) } 96}.$




\textbf{Solution 3}

Analyze that the three-digit integers divisible by $3$ start from $102.$ In the $200$'s, it starts from $201.$ In the $300$'s, it starts from $300.$ We see that the units digits is $0, 1,$ and $2.$

 Write out the $1$- and $2$-digit multiples of $3$ starting from $0, 1,$ and $2.$ Count up the ones that meet the conditions. Then, add up and multiply by $3,$ since there are three sets of three from $1$ to $9.$ Then, subtract the amount that started from $0,$ since the $300$'s ll contain the digit $3.$

 Together, the answer is $3(12+12+12)-12=\boxed{\textbf{(A) } 96}.$




\textbf{Solution 4}

Consider the number of $2$-digit numbers that do not contain the digit $3,$ which is $90-18=72.$ For any of these $2$-digit numbers, we can append $1,5,7,$ or $9$ to reach a desirable $3$-digit number. However, we have $7 \equiv 1\pmod{3},$ and thus we need to count any $2$-digit number $\equiv 2\pmod{3}$ twice. There are $(98-11)/3+1=30$ total such numbers that have remainder $2,$ but $6$ of them $(23,32,35,38,53,83)$ contain $3,$ so the number we want is $30-6=24.$ Therefore, the final answer is $72+24= \boxed{\textbf{(A) } 96}.$




\textbf{Solution 5}

We need to take care of all restrictions. Ranging from $101$ to $999,$ there are $450$ odd $3$-digit numbers. Exactly $\frac{1}{3}$ of these numbers are divisible by $3,$ which is $450\cdot\frac{1}{3}=150.$ Of these $150$ numbers, $\frac{4}{5}$ $\textbf{do not}$ have $3$ in their ones (units) digit, $\frac{9}{10}$ $\textbf{do not}$ have $3$ in their tens digit, and $\frac{8}{9}$ $\textbf{do not}$ have $3$ in their hundreds digit. Thus, the total number of $3$-digit integers is 
 \[900\cdot\frac{1}{2}\cdot\frac{1}{3}\cdot\frac{4}{5}\cdot\frac{9}{10}\cdot\frac{8}{9}=\boxed{\textbf{(A) } 96}.\]




\textbf{Solution 6}

We will start with the numbers that could work. This numbers include \_ \_ $1$, \_ \_ $5$, \_ \_ $7$, \_ \_ $9$. Let's work case by case.

 Case $1$: \_ \_ $1$: The two blanks could be any number that is $2$ mod $3$ that does not include $3$. We have $24$ cases for this case (we could count every case). 

 Case $2$: \_ \_ $5$: The $2$ blanks could be any number that is $1$ mod $3$ that does not include $3$. But we could see that this case has exactly the same solutions to case $1$ because we have a $1-1$ correspondence. We can do the exact same for case $3$. 

 Cases $4$: \_ \_ $9$: We need the blanks to be a multiple of $3$, but does not contain 3. We have $(12, 15, 18, 21, 24, 27, 42, 45, 48, 51, 54, 57, 60, 66, 69, 72, 75, 78, 81, 84, 87, 90, 96, 99)$ which also contains $24$ numbers. Therefore, we have $24 \cdot 4$ which is equal to $\boxed{\textbf{(A) } 96}.$




\textbf{Solution 7}

This problem is solvable by inclusion exclusion principle. There are $\frac{999-105}{6} + 1 = 150$ odd $3$-digit numbers divisible by $3$. We consider the number of $3$-digit numbers divisible by $3$ that contain either $1, 2$ or $3$ digits of $3$.

 For $\underline{AB3}$, $AB$ is any $2$-digit number divisible by $3$, which gives us $\frac{99-12}{3} + 1 = 30$. For $\underline{A3B}$, for each odd $B$, we have $3$ values of $A$ that give a valid case, thus we have $5(3) = 15$ cases. For $\underline{3AB}$, we also have $15$ cases, but when $B=3, 9$, $A$ can equal $0$, so we have $17$ cases. 

 For $\underline{A33}$, we have $3$ cases. For $\underline{3A3}$, we have $4$ cases. For $\underline{33A}$, we have $2$ cases. Finally, there is just one case for $\underline{333}$. 

 By inclusion exclusion principle, we get $150 - 30 - 15 - 17 + 3 + 4 + 2 - 1 = \boxed{\textbf{(A) } 96}$ numbers.

 



\textbf{Solution 8 (only if you don't have time)}

List the numbers that satisfy restriction for $100$ and $200$. Each of them have $12$. Assume that this holds for all other hundreds. Multiply $12$ and $8$ because $300$ doesn't count. The answer is $\boxed{\textbf{(A) } 96}$.

 


	\item (\textit{AMC 12 12B 2018}) The solutions to the equation $(z+6)^8=81$ are connected in the complex plane to form a convex regular polygon, three of whose vertices are labeled $A,B,$ and $C$. What is the least possible area of $\triangle ABC?$
$\textbf{(A) } \frac{1}{6}\sqrt{6} \qquad \textbf{(B) } \frac{3}{2}\sqrt{2}-\frac{3}{2} \qquad \textbf{(C) } 2\sqrt3-3\sqrt2 \qquad \textbf{(D) } \frac{1}{2}\sqrt{2} \qquad \textbf{(E) } \sqrt 3-1$

 \textbf{Answer: }B



\textbf{Solution 1 (Complex Numbers in Polar Form)}

Recall that translations preserve the shapes and the sizes for all objects. We translate the solutions to the given equation $6$ units right, so they become the solutions to the equation $z^8=81.$

 We rewrite $z$ to the polar form 
 \[z=r(\cos\theta+i\sin\theta)=r\operatorname{cis}\theta,\] where $r$ is the magnitude of $z$ such that $r\geq0,$ and $\theta$ is the argument of $z$ such that $0\leq\theta<2\pi.$

 By De Moivre's Theorem, we have 
 \[z^8=r^8\operatorname{cis}(8\theta)={\sqrt3}^8(1),\] from which

 
            \begin{enumerate} 
        \item $r^8={\sqrt3}^8,$ so $r=\sqrt3.$

        \item $\begin{cases} \begin{aligned} \cos(8\theta) &= 1 \\ \sin(8\theta) &= 0 \end{aligned}, \end{cases}$ so $\theta=0,\frac{\pi}{4},\frac{\pi}{2},\frac{3\pi}{4},\pi,\frac{5\pi}{4},\frac{3\pi}{2},\frac{7\pi}{4}.$

            \end{enumerate}
            In the complex plane, the solutions to the equation $z^8=81$ are the vertices of a regular octagon with center $0$ and radius $\sqrt3.$

 The least possible area of $\triangle ABC$ occurs when $A,B,$ and $C$ are the consecutive vertices of the octagon. For simplicity purposes, let $A=\sqrt3\operatorname{cis}\frac{\pi}{4}=\frac{\sqrt6}{2}+\frac{\sqrt6}{2}i, B=\sqrt3\operatorname{cis}\frac{\pi}{2}=\sqrt3i,$ and $C=\sqrt3\operatorname{cis}\frac{3\pi}{4}=-\frac{\sqrt6}{2}+\frac{\sqrt6}{2}i,$ as shown below.
                    \begin{figure}[H]
                        \centering
                        \includegraphics[width=0.5\textwidth]{Solution/AMC_12_12B_2018_16_Solution_Figure_1.png}
                    \end{figure}
                    Note that $\triangle ABC$ has base $AC=\sqrt6$ and height $\sqrt3-\frac{\sqrt6}{2},$ so its area is 
 \[\frac12\cdot\sqrt6\cdot\left(\sqrt3-\frac{\sqrt6}{2}\right)=\boxed{\textbf{(B) } \frac{3}{2}\sqrt{2}-\frac{3}{2}}.\]~MRENTHUSIASM

 



\textbf{Solution 2 (Complex Numbers in Rectangular Form)}

Recall that translations preserve the shapes and the sizes for all objects. We translate the solutions to the given equation $6$ units right, so they become the solutions to the equation $z^8=81.$

 We have
 \begin{align*} z^8 &= 81 \\ z^4 &= 9, -9 \\ z^2 &= 3, -3, 3i, -3i. \\ \end{align*}
Note that:

 
            \begin{enumerate} 
        \item For $z^2=3,$ we get $\boldsymbol{z=\sqrt3,-\sqrt3}.$

        \item For $z^2=-3,$ we get $\boldsymbol{z=\sqrt3i,-\sqrt3i}.$

        \item For $z^2=3i,$ let $z=a+bi$ for some real numbers $a$ and $b.$ We substitute and then expand:
 \begin{align*} (a+bi)^2 &= 3i \\ \left(a^2-b^2\right)+2abi &= 3i. \end{align*}
We equate the real parts and the imaginary parts, respectively, then simplify:
 \begin{align*} a^2 &= b^2, &&(1) \\ ab &= \frac32. &&(2) \end{align*}
Since $ab>0$ in $(2),$ we conclude that $a$ and $b$ must have the same sign. It follows that $(1)$ becomes $a=b.$ By substitution, we get $(a,b)=\left(\frac{\sqrt6}{2},\frac{\sqrt6}{2}\right),\left(-\frac{\sqrt6}{2},-\frac{\sqrt6}{2}\right),$ from which $\boldsymbol{z=\frac{\sqrt6}{2}+\frac{\sqrt6}{2}i,-\frac{\sqrt6}{2}-\frac{\sqrt6}{2}i}.$

        \item For $z^2=-3i,$ we get $\boldsymbol{z=\frac{\sqrt6}{2}-\frac{\sqrt6}{2}i,-\frac{\sqrt6}{2}+\frac{\sqrt6}{2}i}$ by a similar process.

            \end{enumerate}
            We continue with the last two paragraphs of Solution 1 to obtain the answer $\boxed{\textbf{(B) } \frac{3}{2}\sqrt{2}-\frac{3}{2}}.$




\textbf{Solution 3 (Regular Octagon)}

The polygon formed will be a regular octagon since there are $8$ roots of $z^8=81$. By normal math computation, we can figure out that two roots of $z^8=81$ are $\sqrt{3}$ and $-\sqrt{3}$. These will lie on the real axis of the plane. Since it's a regular polygon, there have to be points on the imaginary axis also which will be $\sqrt{3}i$ and $-\sqrt{3}i$.

 Clearly, the rest of the points will lie in each quadrant. The next thing is to get their coordinates (note that to answer this question, we do not need all the coordinates, only 3 consecutive ones are needed).

 The circumcircle of the octagon will have the equation $i^2+r^2=3$. The coordinates of the point in the first quadrant will be equal in magnitude and both positive, so $i=r$. Solving gives $i=r=\frac{\sqrt{3}}{\sqrt{2}}$ (meaning that the root represented is $\frac{\sqrt{3}}{\sqrt{2}}+\frac{\sqrt{3}}{\sqrt{2}}i$).

 This way we can deduce the values of the $8$ roots of the equation to be 
 \[\sqrt{3},-\sqrt{3},-\sqrt{3}i,\sqrt{3}i,\frac{\sqrt{3}}{\sqrt{2}}+\frac{\sqrt{3}}{\sqrt{2}}i,-\frac{\sqrt{3}}{\sqrt{2}}-\frac{\sqrt{3}}{\sqrt{2}}i,\frac{\sqrt{3}}{\sqrt{2}}-\frac{\sqrt{3}}{\sqrt{2}}i,-\frac{\sqrt{3}}{\sqrt{2}}+\frac{\sqrt{3}}{\sqrt{2}}i.\]

 To get the area, $3$ consecutive points such as $\sqrt{3},$ $\frac{\sqrt{3}}{\sqrt{2}}+\frac{\sqrt{3}}{\sqrt{2}}i,$ and $\sqrt{3}i$ can be used. The area can be computed using different methods like using the Shoelace Theorem, or subtracting areas to find the area.The answer you get is $\boxed{\textbf{(B) } \frac{3}{2}\sqrt{2}-\frac{3}{2}}$.

 (This method is not actually as long as it seems if you understand what you're doing while doing it. Also calculations can be made a little easier by solving using $x^8=1$ and multiplying your answer by $\sqrt{3}$).

 



\textbf{Solution 4 (Roots of Unity)}

Now, we need to solve the equation $(z+6)^8 = 81$ where $z = a+bi$. We can substitute this as 
 \[(a+6+bi)^8 = 81\] Now, let $a+6 = q$ for some $q \in \mathbb{Z}$. Thus, the equation becomes $((q+bi)^2)^4 = 81$. Taking it to the other side, we get the equation to be $(q+bi)^2 = 3$. Rearranging variables, we get $(q+bi) = \sqrt{3}$. Plotting this in the complex place, this is a circle centered at the origin and of radius $\sqrt{3}$. 

 The graph of the original equation $(a+6+bi) = \sqrt{3}$ is merely a transformation which doesn't change affect the area. Thus, we can find the minimum area of the transformed equation $(q+bi)^2 = 3$. Using Roots of Unity, we know that the roots of the equation lie at $0, \frac{\pi}{4}, \frac{\pi}{2}, \frac{3\pi}{4}, \ldots, 2\pi$ radians from the origin. 

 We can quickly notice that the area of the roots will be smallest with points at $\frac{\pi}{4}, \frac{\pi}{2}, \frac{3\pi}{4}$. Using trigonometry, we get the respective roots to be $(\operatorname{Re}(z), \operatorname{Im}(z)) \in \left\{\left(\sqrt{3},0\right), \left(\frac{\sqrt{6}}{2}, \frac{\sqrt{6}}{2}\right), \left(0,\sqrt{3}\right)\right\}$. Using the Shoelace Theorem, the area quickly comes out to be $\frac{3\sqrt{2}-3}{2} = \boxed{\textbf{(B) } \frac{3}{2}\sqrt{2}-\frac{3}{2}}.$




\textbf{Solution 5 (Trigonometry)}

It is easy to see that the figure the solutions form is that of an regular octagon $ABCDEFGH$ With center $O$ such that $AO = \sqrt{3}.$ Also, the minimum area of a triangle formed by the solutions would be that of Triangle $ABC$, since the three points are consecutive vertices of the octagon. By dropping an altitude from $O$ to $AB$, note that $\frac{1}{2} AB = \sqrt{3} \sin(22.5).$ Thus, $AB = 2\sqrt{3} \sin(22.5).$ Now given the length of $AB$ we need to find the area of $ABC,$ which, by simple trigonometry, is given as 
 \[AB \cos(22.5) \cdot AB \sin(22.5) = 12 \sin^3(22.5) \cdot \cos(22.5).\] By the half angle formulas, $\sin(22.5) = \sqrt{\frac{1 - \sin(45)}{2}} = \frac{\sqrt{2-\sqrt{2}}}{2},$ and $\cos(22.5)$ is its conjugate $\frac{\sqrt{2+\sqrt{2}}}{2}.$ Doing some expanding and simplifying, 
 \[12 \cdot \left ( \frac{\sqrt{2-\sqrt{2}}}{2} \right)^3 \cdot \left (\frac{\sqrt{2+\sqrt{2}}}{2} \right) = \frac{3\sqrt{12-8\sqrt{2}}}{4}.\] Further simplifying, we obtain that the desired expression is $\frac{3 \sqrt{3 - 2\sqrt{2}}}{2}.$ Note that $(\sqrt{2} - 1)^2 = 3 - 2\sqrt{2},$ so our answer comes out to be $\frac{3(\sqrt{2}-1)}{2} = \boxed{\textbf{(B) } \frac{3}{2}\sqrt{2}-\frac{3}{2}}.$



	\item (\textit{AMC 12 12B 2018}) Let $p$ and $q$ be positive integers such that 
 \[\frac{5}{9} < \frac{p}{q} < \frac{4}{7}\]and $q$ is as small as possible. What is $q-p$?


 $\textbf{(A) } 7 \qquad \textbf{(B) } 11 \qquad \textbf{(C) } 13 \qquad \textbf{(D) } 17 \qquad \textbf{(E) } 19$

 \textbf{Answer: }A



\textbf{Solution 1 (Generalization)}

More generally, let $a,b,c,d,p,$ and $q$ be positive integers such that $bc-ad=1$ and 
 \[\frac ab < \frac pq < \frac cd.\]From $\frac ab < \frac pq,$ we have $bp-aq>0,$ or 
 \[bp-aq\geq1. \hspace{15mm} (1)\]From $\frac pq < \frac cd,$ we have $cq-dp>0,$ or 
 \[cq-dp\geq1. \hspace{15mm} (2)\]Since $bc-ad=1,$ note that:

 
            \begin{enumerate} 
        \item Multiplying $(1)$ by $d,$ multiplying $(2)$ by $b,$ and adding the results, we get $q\geq b+d.$

        \item Multiplying $(1)$ by $c,$ multiplying $(2)$ by $a,$ and adding the results, we get $p\geq a+c.$

            \end{enumerate}
            To minimize $q,$ we set $q=b+d,$ from which $p=a+c.$ Together, we can prove that 
 \[\frac{a+c-1}{b+d}\leq\frac ab<\frac{a+c}{b+d}<\frac cd\leq\frac{a+c+1}{b+d}. \hspace{15mm} (\bigstar)\]For this problem, we have $a=5,b=9,c=4,$ and $d=7,$ so $p=a+c=9$ and $q=b+d=16.$ The answer is $q-p=\boxed{\textbf{(A) } 7}.$
<u>\textbf{Remark}</u>

 We will prove each part of the compound inequality in $(\bigstar):$

            \begin{enumerate} 
        \item $\frac{a+c-1}{b+d}\leq\frac ab$ and $\frac cd\leq\frac{a+c+1}{b+d}$Let $\frac ab=k,$ so $a=bk.$ The precondition $bc-ad=1$ becomes $bc-bdk=1,$ so $c-dk=\frac1b.$ 

 It follows that 
 \[\frac{a+c-1}{b+d}=\frac{bk+c-1}{b+d}=k+\frac{c-dk-1}{b+d}=k+\frac{\frac1b-1}{b+d}=k+\frac{1-b}{b(b+d)}\leq k.\] Moreover, the equality case occurs if and only if $b=1.$ 

 We can prove $\frac cd\leq\frac{a+c+1}{b+d}$ by the same process. Similarly, the equality case occurs if and only if $d=1.$

        \item $\frac ab<\frac{a+c}{b+d}<\frac cd$Let $\frac ab=k_1$ and $\frac cd=k_2,$ so $a=bk_1$ and $c=dk_2.$ 

 It follows that
 \begin{alignat*}{10} \frac{a+c}{b+d}&=\frac{bk\_1+dk\_2}{b+d}&&=k\_1+\frac{dk\_2-dk\_1}{b+d}&&=k\_1+\frac{d(k\_2-k\_1)}{b+d}&&>k\_1, \\ \frac{a+c}{b+d}&=\frac{bk\_1+dk\_2}{b+d}&&=k\_2+\frac{bk\_1-bk\_2}{b+d}&&=k\_2+\frac{b(k\_1-k\_2)}{b+d}&&<k\_2. \end{alignat*}Moreover, this part of $(\bigstar)$ is independent of the precondition $bc-ad=1.$

            \end{enumerate}
            



\textbf{Solution 2 (Generalization)}

Define variables $a,b,c,d,p,$ and $q$ as Solution 1 does. Moreover, this solution refers to inequalities $(1)$ and $(2)$ in Solution 1.

 Note that
 \begin{align*} \frac{1}{bd}&=\frac{bc-ad}{bd} \\ &=\frac cd - \frac ab \\ &=\left(\frac cd - \frac pq\right)+\left(\frac pq - \frac ab\right) \\ &=\frac{cq-dp}{dq}+\frac{bp-aq}{bp} \\ &\geq\frac{1}{dq}+\frac{1}{bq}. \end{align*}
Multiplying both sides of $\frac{1}{bd}\geq\frac{1}{dq}+\frac{1}{bq}$ by $bdq,$ we get $q\geq b+d.$

 For this problem, we have $a=5,b=9,c=4,$ and $d=7,$ so $q\geq b+d=16.$ At $q=16,$ we have 
 \[\frac{8}{16}<\frac59<\frac{9}{16}<\frac47<\frac{10}{16},\]from which $p=9.$ Therefore, the answer is $q-p=\boxed{\textbf{(A) } 7}.$

 Alternatively, refer to the <u>\textbf{Remark}</u> section in Solution 1 for further generalizations. 

 



\textbf{Solution 3 (Substitution)}

Inverting the given inequality we get 
 \[\frac{7}{4} < \frac{q}{p} < \frac{9}{5},\]which simplifies to 
 \[35p < 20q < 36p.\]We can now substitute $q = p + k.$ Note we need to find $k:$ 
 \[35p < 20p + 20k < 36p,\]which simplifies to 
 \[15p < 20k < 16p.\]Clearly since $20k < 16p$, $p>k.$ We will now substitute $p = k + x$ to get 
 \[15k + 15x < 20k < 16k + 16x.\]The inequality $15k + 15x < 20k$ simplifies to $3x < k.$

 The inequality $20k < 16k + 16x$ simplifies to $k < 4x.$

 Combining the two inequalities, we get 
 \[3x < k < 4x.\]Since $x$ and $k$ are integers, the smallest values of $x$ and $k$ that satisfy the above equation are $2$ and $7$ respectively. Substituting these back in, we arrive with an answer of $\boxed{\textbf{(A) } 7}.$




\textbf{Solution 4 (Substitution)}

Because $q$ and $p$ are positive integers with $p<q,$ we can let $q=p+k$ where $k\in{\mathbb{Z}}.$ Now, the problem condition reduces to 
 \[\frac{5}{9}<\frac{p}{p+k}<\frac{4}{7}.\]

 Our first inequality is $\frac{5}{9}<\frac{p}{p+k},$ which gives us $5p+5k<9p,$ or $\frac{5}{4}k<p.$

 Our second inequality is $\frac{p}{p+k}<\frac{4}{7},$ which gives us $7p<4p+4k,$ or $p<\frac{4}{3}k.$

 Hence, we have $\frac{5}{4}k<p<\frac{4}{3}k,$ or $15k<12p<16k.$

 It is clear that we are aiming to find the least positive integer value of $k$ such that there is at least one value of $p$ that satisfies the inequality.

 Now, simple casework through the answer choices of the problem reveals that $q-p=p+k-p=k\ge{\boxed{\textbf{(A) } 7}}.$




\textbf{Solution 5 (Manipulation)}

Subtract $\frac{1}{2}$ from both sides of the equation. This implies that 
 \[\frac{1}{18} < \frac{p}{q}-\frac{1}{2} < \frac{1}{14}.\] Then for $q$ to be minimized, $\frac{p}{q}-\frac{1}{2}$ must be $\frac{1}{16},$ so $\frac{p}{q}$ is $\frac{9}{16}$ and $q-p$ is $\boxed{\textbf{(A) } 7}.$




\textbf{Solution 6 (Solves for p)}

Cross-multiply the inequality to get 
 \[35q < 63p < 36q.\]Then, we have $0 < 63p-35q < q,$ or 
 \[0 < 7(9p-5q) < q.\]Since $p$ and $q$ are integers, $9p-5q$ is an integer. To minimize $q,$ start from $9p-5q=1,$ which gives $p=\frac{5q+1}{9}.$ This limits $q$ to be greater than $7,$ so test values of $q$ starting from $q=8.$ However, $q=8$ to $q=14$ do not give integer values of $p.$ 

 Once $q>14,$ it is possible for $9p-5q$ to be equal to $2,$ so $p$ could also be equal to $\frac{5q+2}{9}.$ The next value, $q=15,$ is not a solution, but $q=16$ gives $p=\frac{5\cdot 16 + 1}{9} = 9.$ Thus, the smallest possible value of $q$ is $16,$ and the answer is $16-9= \boxed{\textbf{(A) } 7}.$




\textbf{Solution 7 (Inspection)}

Start with $\frac{5}{9}.$ Repeat the following process until you arrive at the answer: if the fraction is less than or equal to $\frac{5}{9},$ add $1$ to the numerator; otherwise, if it is greater than or equal to $\frac{4}{7},$ add one $1$ to the denominator. We have 
 \[\frac{5}{9}, \frac{6}{9}, \frac{6}{10}, \frac{6}{11}, \frac{7}{11}, \frac{7}{12}, \frac{7}{13}, \frac{8}{13}, \frac{8}{14}, \frac{8}{15}, \frac{9}{15}, \frac{9}{16}.\]

 Therefore, the answer is $16 - 9 = \boxed{\textbf{(A) } 7}.$




\textbf{Solution 8 (Inspection)}

Checking possible fractions within the interval can get us to the answer, but only if we do it with more skill.The interval can also be written as $0.5556<x<0.5714.$ This represents fraction with the numerator a little bit more than half the denominator. Every fraction we consider must not exceed this range. 

 The denominators to be considered are $9,10,11,12,\ldots.$ We check $\frac{6}{10}, \frac{6}{11}, \frac{7}{12}, \frac{7}{13}, \frac{8}{15}, \frac{9}{16}.$ At this point we know that we've got our fraction and our answer is $16-9=\boxed{\textbf{(A) } 7}.$

 The inspection was made faster by considering the fact that $\frac{a+1}{b+1}>\frac{a}{b}.$

 So, once a fraction was gotten which was greater than $\frac{4}{7}$ we jump to the next denominator. 

 We then make sure we consider fractions with higher positive difference between the denominator and numerator. And we also do not forget that the numerator must be greater than half of the denominator.

 ($\frac{8}{14}$ was obviously skipped because it is equal to $\frac{4}{7}.$)

 



\textbf{Solution 9 (Graphing)}

Graph the regions $y > \frac{5}{9}x$ and $y < \frac{4}{7}x.$ Note that the lattice point $(16,9)$ is the smallest magnitude one which appears within the region bounded by the two graphs. Thus, our fraction is $\frac{9}{16}$ and the answer is $16-9= \boxed{\textbf{(A) } 7}.$

 Remark: This also gives an intuitive geometric proof of the mediant using vectors.

 



\textbf{Solution 10 (Answer Choices)}

As the other solutions do, the mediant $\frac{9}{16}$ is between the two fractions, with a difference of $\boxed{\textbf{(A) } 7}.$ Suppose that the answer was not $\textbf{(A)},$ then the answer must be $\textbf{(B)}$ or $\textbf{(C)}$ as otherwise $p$ would be negative. Then, the possible fractions with lower denominator would be $\frac{k-11}{k}$ for $k=12,13,14,15$ and $\frac{k-13}{k}$ for $k=14,15,$ which are clearly not anywhere close to $\frac{4}{7}\approx 0.6.$




\textbf{Solution 11 (Answer Choices)}

In ascending order, we can use answer choices, values for $q-p,$ as a method of figuring out our answer through the means of substitution. Let the assumed difference be $7.$ Then, $p=q-7.$ We thus have two inequalities: $\frac{5}{9} < \frac{q-7}{q}$ and $\frac{q-7}{q} < \frac{4}{7}.$

 Solving for $q$ in these equalities, we get $\frac{63}{4} < q < \frac{49}{3}.$ So, $q$ is between $15.75$ and $16.\overline{3},$ making it $16$ as $q$ is a positive integer (again, at this point, this is still an assumption). This would set $p=16-7=9.$

 Since $\frac{5}{9} < \frac{9}{16} < \frac{4}{7},$ the minimum difference is $\boxed{\textbf{(A) } 7}.$




\textbf{Solution 12 (Educated Guess)}

Assume that the difference $\frac{p}{q} - \frac{5}{9}$ results in a fraction of the form $\frac{1}{9q}.$ Then, 
 \[9p - 5q = 1.\]Also assume that the difference $\frac{4}{7} - \frac{p}{q}$ results in a fraction of the form $\frac{1}{7q}.$ Then, 
 \[4q - 7p = 1.\]Solving the system of equations yields $q=16$ and $p=9.$ Therefore, the answer is $\boxed{\textbf{(A) } 7}.$

 Refer to Solution 1 for the full justification.

 



\textbf{Solution 13}

Notice the following property.
 \[\begin{cases}     \frac{a}{b} > 1 &\Rightarrow \frac{a}{b} > \frac{a + 1}{b + 1} > \frac{a + 2}{b + 2} > \cdots \\     \frac{a}{b} < 1 &\Rightarrow \frac{a}{b} < \frac{a + 1}{b + 1} < \frac{a + 2}{b + 2} < \cdots \end{cases}\]Using the property above, the following fact could be derived.
 \[\frac{35}{63} < \frac{36}{64} < \cdots < \frac{36}{63} \qquad \left( \because \lim_{x \rightarrow \infty} \frac{35 + x}{63 + x} = 1 \right)\]First, the range for $x$ could be found:
 \begin{align*}     \frac{35 + x}{63 + x} &= \frac{36}{63} \\     35 \cdot 63 + 63x &= 36 \cdot 63 + 36x \\     27x &= 63 \\     x &= \frac{7}{3}. \end{align*}
In other words, the only possible value for $x$ are $1$ and $2$. It is evident that $q$ will be smallest when $x = 1$. Therefore, $q - p = 16 - 9 = \boxed{\textbf{(A) } 7}$.

 



\textbf{Solution 14 (Trial and Error)}

If $\frac{5}{9} < \frac{p}{q} < \frac{4}{7}$, then $\frac{7}{4} < \frac{q}{p} < \frac{9}{5}$.Since we are given easy values, we can easily find the decimal value for each fraction and minimize $p$ as much as possible.

 We now have: 
 \[\frac{7}{4} = 1.75, \frac{9}{5} = 1.8.\]We obtain $1.75p < q < 1.8p$. Since $p=10$ doesn't work, we try for $p<10$.

 After only one try, we see $p=9$ works, means $q=16.$ The final answer for $q-p=\boxed{\textbf{(A) } 7}$.

 Note: INMO 2005 and IOQM 2023 have a very similar problem. However, the numbers are much more difficult to approximate, which means the above solutions have a much better chance of working ;)

 


 


	\item (\textit{AMC 12 12B 2018}) A function $f$ is defined recursively by $f(1)=f(2)=1$ and 
 \[f(n)=f(n-1)-f(n-2)+n\]for all integers $n \geq 3$. What is $f(2018)$?


 $\textbf{(A) } 2016 \qquad \textbf{(B) } 2017 \qquad \textbf{(C) } 2018 \qquad \textbf{(D) } 2019 \qquad \textbf{(E) } 2020$

 \textbf{Answer: }B



\textbf{Solution 1 (Algebra)}

For all integers $n \geq 7,$ note that
 \begin{align*} f(n)&=f(n-1)-f(n-2)+n \\ &=[f(n-2)-f(n-3)+n-1]-f(n-2)+n \\ &=-f(n-3)+2n-1 \\ &=-[f(n-4)-f(n-5)+n-3]+2n-1 \\ &=-f(n-4)+f(n-5)+n+2 \\ &=-[f(n-5)-f(n-6)+n-4]+f(n-5)+n+2 \\ &=f(n-6)+6. \end{align*}
It follows that
 \begin{align*} f(2018)&=f(2012)+6 \\ &=f(2006)+12 \\ &=f(2000)+18 \\ & \ \vdots \\ &=f(2018-6\cdot336)+6\cdot336 \\ &=f(2)+2016 \\ &=\boxed{\textbf{(B) } 2017}. \end{align*}
~MRENTHUSIASM

 



\textbf{Solution 2 (Algebra)}

For all integers $n\geq3,$ we rearrange the given equation: 
 \[f(n)-f(n-1)+f(n-2)=n. \hspace{28.25mm}(1)\]For all integers $n\geq4,$ it follows that 
 \[f(n-1)-f(n-2)+f(n-3)=n-1. \hspace{15mm}(2)\]For all integers $n\geq4,$ we add $(1)$ and $(2):$ 
 \[f(n)+f(n-3)=2n-1. \hspace{38.625mm}(3)\]For all integers $n\geq7,$ it follows that 
 \[f(n-3)+f(n-6)=2n-7. \hspace{32mm}(4)\]For all integers $n\geq7,$ we subtract $(4)$ from $(3):$ 
 \[f(n)-f(n-6)=6. \hspace{47.5mm}(5)\]From $(5),$ we have the following system of $\frac{2018-8}{6}+1=336$ equations:
 \begin{align*} f(2018)-f(2012)&=6, \\ f(2012)-f(2006)&=6, \\ f(2006)-f(2000)&=6, \\ & \ \vdots \\ f(8)-f(2)&=6. \end{align*}
We add these equations up to get 
 \[f(2018)-f(2)=6\cdot336=2016,\] from which $f(2018)=f(2)+2016=\boxed{\textbf{(B) } 2017}.$




\textbf{Solution 3 (Finite Differences)}

Preamble: In this solution, we define the sequence $A$ to satisfy $a_n = f(n),$ where $a_n$ represents the $n$th term of the sequence $A.$ This solution will show a few different perspectives. Even though it may not be as quick as some of the solutions above, I feel like it is an interesting concept, and may be more motivated. 

 To begin, we consider the sequence $B$ formed when we take the difference of consecutive terms between $A.$ Define $b_n = a_{n+1} - a_n.$ Notice that for $n \ge 4,$ we have 

 $\begin{cases}\begin{aligned} a_{n+1} &= a_{n} - a_{n-1} + (n+1) \\ a_{n} &= a_{n-1} - a_{n-2} + n \end{aligned}.\end{cases}$

  Notice that subtracting the second equation from the first, we see that $b_{n} = b_{n-1} - b_{n-2} + 1.$ 

 If you didn’t notice that $B$ repeated directly in the solution above, you could also, possibly more naturally, take the finite differences of the sequence $b_n$ so that you could define $c_n = b_{n+1} - b_n.$ Using a similar method as above through reindexing and then subtracting, you could find that $c_n = c_{n-1} - c_{n-2}.$ The sum of any six consecutive terms of a sequence which satisfies such a recursion is $0,$ in which you have that $b_{n} = b_{n+6}.$ In the case in which finite differences didn’t reduce to such a special recursion, you could still find the first few terms of $C$ to see if there are any patterns, now that you have a much simpler sequence. Doing so in this case, it can also be seen by seeing that the sequence $C$ looks like 
 \[\underbrace{2, 1, -1, -2, -1, 1,}_{\text{cycle period}} 2, 1, -1, -2, -1, 1, \ldots\] in which the same result follows. 

 Using the fact that $B$ repeats every six terms, this motivates us to look at the sequence $B$ more carefully. Doing so, we see that $B$ looks like 
 \[\underbrace{2, 3, 2, 0, -1, 0,}_{\text{cycle period}} 2, 3, 2, 0, -1, 0, \ldots\] (If you tried pattern finding on sequence $B$ directly, you could also arrive at this result, although I figured defining a second sequence based on finite differences was more motivated.)

 Now, there are two ways to finish. 

 Finish Method \#1: Notice that any six consecutive terms of $B$ sum to $6,$ after which we see that $a_n = a_{n-6} + 6.$ Therefore, $a_{2018} = a_{2012} + 6 = \cdots = a_{2} + 2016 = \boxed{\textbf{(B) } 2017}.$

 Finish Method \#2: Notice that $a_{2018} = a_{2017} - a_{2016} + 2018 = B_{2016} + 2018 = -1 + 2018 = \boxed{\textbf{(B) } 2017}.$




\textbf{Solution 4 (Pattern)}

Start out by listing some terms of the sequence. 
 \begin{align*} f(1)&=1 \\ f(2)&=1 \\ f(3)&=3 \\ f(4)&=6 \\ f(5)&=8 \\ f(6)&=8 \\  f(7)&=7 \\ f(8)&=7 \\ f(9)&=9 \\ f(10)&=12 \\ f(11)&=14 \\ f(12)&=14 \\ f(13)&=13 \\ f(14)&=13 \\ f(15)&=15 \\ & \ \vdots \end{align*}
Notice that $f(n)=n$ whenever $n$ is an odd multiple of $3$, and the pattern of numbers that follow will always be $+3$, $+2$, $+0$, $-1$, $+0$, $+2$.The largest odd multiple of $3$ smaller than $2018$ is $2013$, so we have
 \begin{align*} f(2013)&=2013 \\ f(2014)&=2016 \\ f(2015)&=2018 \\ f(2016)&=2018 \\ f(2017)&=2017 \\ f(2018)&=\boxed{\textbf{(B) } 2017}. \end{align*}





\textbf{Solution 5 (Pattern)}

Writing out the first few values, we get
 \[1,1,3,6,8,8,7,7,9,12,14,14,13,13,15,18,20,20,19,19,\ldots.\] We see that every number $x$ where $x \equiv 1\pmod 6$ has $f(x)=x,f(x+1)=f(x)=x,$ and $f(x-1)=f(x-2)=x+1.$ The greatest number that's $1\pmod{6}$ and less than $2018$ is $2017,$ so we have $f(2017)=f(2018)=\boxed{\textbf{(B) } 2017}.$




\textbf{Solution 6 (Complex Numbers)}


 \begin{align*} f(n)&=f(n-1)-f(n-2)+n \\ f(n-1)&=f(n-2)-f(n-3)+n-1 \end{align*}
Subtracting those two and rearranging gives 
 \begin{align*} f(n)-2f(n-1)+2f(n-2)-f(n-3)&=1 \\ f(n-1)-2f(n-2)+2f(n-3)-f(n-4)&=1 \end{align*}
Subtracting those two gives $f(n)-3f(n-1)+4f(n-2)-3f(n-3)+f(n-4)=0.$

 The characteristic polynomial is $x^4-3x^3+4x^2-3x+1=0.$
$x=1$ is a root, so using synthetic division results in $(x-1)(x^3-2x^2+2x-1)=0.$
$x=1$ is a root, so using synthetic division results in $(x-1)^2(x^2-x+1)=0.$
$x^2-x+1=0$ has roots $x=\frac{1}{2}\pm\frac{i\sqrt{3}}{2}.$

 And 
 \[f(n)=(An+D)\cdot1^n+B\cdot\left(\frac{1}{2}-\frac{i\sqrt{3}}{2}\right)^n+C\cdot\left(\frac{1}{2}+\frac{i\sqrt{3}}{2}\right)^n.\]Plugging in $n=1$, $n=2$, $n=3$, and $n=4$ results in a system of $4$ linear equations$\newline$Solving them gives $A=1, \ B=\frac{1}{2}-\frac{i\sqrt{3}}{2}, \ C=\frac{1}{2}+\frac{i\sqrt{3}}{2}, \ D=1.$ Note that you can guess $A=1$ by answer choices.

 So plugging in $n=2018$ results in 
 \begin{align*} 2018+1+\left(\frac{1}{2}-\frac{i\sqrt{3}}{2}\right)^{2019}+\left(\frac{1}{2}+\frac{i\sqrt{3}}{2}\right)^{2019}&=2019+(\cos(-60^{\circ})+i\sin(-60^{\circ}))^{2019})+(\cos(60^{\circ})+i\sin(60^{\circ}))^{2019}) \\ &=2019+(\cos(-60^{\circ}\cdot2019)+i\sin(-60^{\circ}\cdot2019))+(\cos(60^{\circ}\cdot2019)+i\sin(60^{\circ}\cdot2019)) \\ &=\boxed{\textbf{(B) } 2017}. \end{align*}
~ryanbear

 



\textbf{Solution 7}

We utilize patterns to solve this equation:
 \begin{align*} f(3)&=3, \\ f(4)&=6, \\ f(5)&=8, \\ f(6)&=8, \\ f(7)&=7, \\ f(8)&=7. \end{align*}
We realize that the pattern repeats itself. For every six terms, there will be four terms that we repeat, and two terms that we don't repeat. We will exclude the first two for now, because they don't follow this pattern. 

 First, we need to know whether or not $2016$ is part of the skip or repeat. We notice that $f(6),f(12), \ldots,f(6n)$ all satisfy $6+6(n-1)=n,$ and we know that $2016$ satisfies this, leaving $n=336.$ Therefore, we know that $2016$ is part of the repeat section. But what number does it repeat? 

 We know that the repeat period is $2,$ and it follows that pattern of $1,1,8,8,7,7.$ Again, since $f(6) = f(5)$ and so on  for the repeat section, $f(2016)=f(2015),$ so we don't need to worry about which one, since it repeats with period $2.$ We see that the repeat pattern of $f(6),f(12),\ldots,f(6n)$ follows $8,14,20,$ it is an arithmetic sequence with common difference $6.$ Therefore, $2016$ is the $335$th term of this, but including $1,$ it is $336\cdot6+1=\boxed{\textbf{(B) } 2017}.$

 
Edited by TheBinaryBlackHole

 



\textbf{Solution 8}

We list out the first $13$ terms and their differences: 

 
                    \begin{figure}[H]
                        \centering
                        \includegraphics[width=0.5\textwidth]{Solution/AMC_12_12B_2018_18_Solution_Figure_1.png}
                    \end{figure}
                     We see the differences are clearly repeating every $6$ terms here; we see that $2018\equiv2\pmod6$. We see that the values for $2$ and $8$ are both one less, so we conclude our answer is $2018-1=\boxed{\textbf{(B)}~2017}.$



	\item (\textit{AMC 12 12B 2018}) Mary chose an even $4$-digit number $n$. She wrote down all the divisors of $n$ in increasing order from left to right: $1,2,\ldots,\dfrac{n}{2},n$. At some moment Mary wrote $323$ as a divisor of $n$. What is the smallest possible value of the next divisor written to the right of $323$?


 $\textbf{(A) } 324 \qquad \textbf{(B) } 330 \qquad \textbf{(C) } 340 \qquad \textbf{(D) } 361 \qquad \textbf{(E) } 646$

 \textbf{Answer: }C



\textbf{Solution 1 (Inequalities)}

Let $d$ be the next divisor written to the right of $323.$

 If $\gcd(323,d)=1,$ then 
 \[n\geq323d>323^2>100^2=10000,\] which contradicts the precondition that $n$ is a $4$-digit number.

 It follows that $\gcd(323,d)>1.$ Since $323=17\cdot19,$ the smallest possible value of $d$ is $17\cdot20=\boxed{\textbf{(C) } 340},$ from which 
 \[n=\operatorname{lcm}(323,d)=17\cdot19\cdot20=6460.\]~MRENTHUSIASM ~tdeng

 



\textbf{Solution 2 (Inequalities)}

Let $d$ be the next divisor written to the right of $323.$

 Since $n$ is even and $323=17\cdot19,$ we have $n=2\cdot17\cdot19\cdot k=646k$ for some positive integer $k.$ Moreover, since $1000\leq n\leq9998,$ we get $2\leq k\leq15.$ As $d>323,$ it is clear that $d$ must be divisible by $17$ or $19$ or both.

 Therefore, the smallest possible value of $d$ is $17\cdot20=\boxed{\textbf{(C) } 340},$ from which 
 \[n=\operatorname{lcm}(323,d)=17\cdot19\cdot20=6460.\]~MRENTHUSIASM ~bjhhar

 



\textbf{Solution 3 (Quick)}

The prime factorization of $323$ is $17 \cdot 19$. Our answer must be a multiple of either $17$ or $19$ or both. Since $17 < 19$, the next smallest divisor that is divisble by $17$ would be $323 + 17 = \boxed{\textbf{(C) } 340}$




\textbf{Solution 4 (Answer Choices)}

Since prime factorizing $323$ gives you $17 \cdot 19$, the desired answer needs to be a multiple of $17$ or $19$, this is because if it is not a multiple of $17$ or $19$, $n$ will be more than a $4$ digit number. For example, if the answer were to instead be $324$, $n$ would have to be a multiple of $2^2\cdot3^4\cdot17\cdot19$ for both $323$ and $324$ to be a valid factor, meaning $n$ would have to be at least $104652$, which is too big. Looking at the answer choices, $\textbf{(A)}$ and $\textbf{(B)}$ are both not a multiple of neither $17$ nor $19$, $\textbf{(C)}$ is divisible by $17$. $\textbf{(D)}$ is divisible by $19$, and $\textbf{(E)}$ is divisible by both $17$ and $19$. Since $\boxed{\textbf{(C) } 340}$ is the smallest number divisible by either $17$ or $19$ it is the answer. Checking, we can see that $n$ would be $6460$, a $4$-digit number. Note that $n$ is also divisible by $2$, one of the listed divisors of $n$. (If $n$ was not divisible by $2$, we would need to look for a different divisor.)

 -Edited by Mathandski

 



\textbf{Solution 5 (Answer Choices)}

Note that $323$ multiplied by any of the answer choices results in a $5$ or $6$-digit $n$. So, we need a choice that shares a factor(s) with $323$, such that the factors we'll need to add to the prime factorization of $n$ (in result to adding the chosen divisor) won't cause our number to multiply to more than $4$ digits. The prime factorization of $323$ is $17\cdot19$, and since we know $n$ is even, our answer needs to be 

 
            \begin{itemize}
        \item even
            \end{itemize}
            
            \begin{itemize}
        \item has a factor of $17$ or $19$
            \end{itemize}
            We see $340$ achieves this and is the smallest to do so ($646$ being the other). So, we get $\boxed{\textbf{(C) } 340}$.

 



\textbf{Solution 6 (Olympiad Proof )}

This is not the fastest solution, but if i saw this question on an Olympiad/AIME, where there are no answer choices, and my work counted, this is what i would do (for the purpose of this question, a=323, b=answer, c=the four digit number);

 $\newline$Note that for any $a,b,c \in \mathbb{Z}^+$, if $a|c$ and $b|c$ then lcm$(a,b)|c$ (this is because if something divides a or b, it must divide c, and thus the max of all of the prime factors, ie: lcm, divides c) since $a=323,b > a$, if $gcd(a,b)=1$ then $lcm(a,b)=ab$ and thus $ab|c \implies c>a^2 \implies c>100^2 \implies$ c is not a four digit number.$\newline$ thus, $gcd(a,b)\neq1$. This implies that either $17|b$, or $19|b$, or both.

 $\newline$ Case 1: $17|b$, $19\not|b$. We let $b=17b'$, and by Euclid's Lemma, $19\not|b'$. Then, $lcm(323,b)|c \implies 17(lcm(19,b'))|c$. Since we already established that, $19\not|b$ (and since 19 is prime, if it does not divide a number it is coprime to that number), $17*19*b'|c \implies 323b'|c$. Since $b=17b'>19*17$, $b' \geq 20$. A quick check shows $b'=20, b=340$ suffices.$\newline$ Now, let us show that there are no such numbers less than 340.$\newline$ Presume there exists such a number, $n \in \mathbb{Z}^+$ is in the range $(323,340)$. By hypothesis, there is a $d>1$ such that $d|323$, $d|n$. By properties of divisibility $d|n-323$. the maximum possible value of $n-323$ is $17^-$ (basically an arbitrary amount smaller than 17). But, since $d>1$ and $d|323, d \in \{17,19,323\}$. Of which, the minimum value is d=17. but, $17>17^-$ so there is no such d, and no such n.$\newline$ Thus, our answer is just $\boxed{\textbf{(C) } 340}$.$\newline$ ~Stereotypicalmathnerd

 



\textbf{Solution 7}

Let the value of the next divisor be $x$. Then it follows that $lcm(323,x)|n$. For the sake of this problem, suppose $lcm(323,x)=n$, since it seems unlikely that $lcm(323,x)\leq 4999$. From here, there are two ways to solve this problem (note that $323=17\times19)$. $\newline$
$\textbf{Ending 1.1 }$ 

 You can simply prime factorize each answer choice and find one that will produce a valid value for $n$. $\newline$
$\textbf{Ending 1.2 }$ 

 Alternatively, if you know the identity relating $gcd$ and $lcm$ ($gcd(a,b)\times lcm(a,b)=a\times b$), this problem can be rewritten as $gcd(323,x)\times lcm(323,x)=323\times x$. Since $lcm(323,x)=n$, this means $n=\frac{323\times x}{gcd(323,x)}$.

 Using either ending will produce the final answer of $\boxed{\textbf{(C) } 340}$ (This means $n=\frac{323\times 340}{gcd(323,340)} \to n=\frac{17^2 \times 19 \times 20}{17} \to n=6460$, which checks out).$\newline$




\textbf{Solution 8}

The key point is that $n$ is even. Also, $323n < 10000 \implies n \leq 30$. The key idea is that $323 = 17 \cdot 19$. If we were to divide by some $n < 17$, then dividing by 17 yields a smaller value, which is not allowed. Hence, consider $n = 18$ in which $19 \cdot 18 = 342$ is obtained. 

 Now we check the next value, when $n = 20$ so we can check 19. In that case, $17 \cdot 20 = 340$ is formed.

 We now prove no other values exist. Simply speaking, if $n > 20$, then consider $323 \cdot n$. If $n$ is a prime, dividing by $n$ gives 323, and the next number is anyways $17 \cdot n$. This however is minimized at $n = 20$, hence no cases arise. Similarly, if $n$ is composite with a prime factorization, one can still rewrite it of the form $323 \cdot n$. But all prime factors of $n$ are less than or equal to 15, hence the minimum is still at when we divide by 19, in which $17 \cdot n$. But this is still minimized at $n = 20$, hence no cases arise. $\blacksquare$

 The answer is $\boxed{340}$, and we are done. $\square$



	\item (\textit{AMC 12 12B 2018}) Let $ABCDEF$ be a regular hexagon with side length $1$. Denote by $X$, $Y$, and $Z$ the midpoints of sides $\overline {AB}$, $\overline{CD}$, and $\overline{EF}$, respectively. What is the area of the convex hexagon whose interior is the intersection of the interiors of $\triangle ACE$ and $\triangle XYZ$?


 $\textbf{(A)}\ \frac {3}{8}\sqrt{3} \qquad \textbf{(B)}\ \frac {7}{16}\sqrt{3} \qquad \textbf{(C)}\ \frac {15}{32}\sqrt{3} \qquad  \textbf{(D)}\ \frac {1}{2}\sqrt{3} \qquad \textbf{(E)}\ \frac {9}{16}\sqrt{3}$

 \textbf{Answer: }C



\textbf{Solution 1 (Area Addition)}


                    \begin{figure}[H]
                        \centering
                        \includegraphics[width=0.5\textwidth]{Solution/AMC_12_12B_2018_20_Solution_Figure_1.png}
                    \end{figure}
                    

 The desired area (hexagon $MPNQOR$) consists of an equilateral triangle ($\triangle MNO$) and three right triangles ($\triangle MPN,\triangle NQO,$ and $\triangle ORM$).

 Notice that $\overline {AD}$ (not shown) and $\overline {BC}$ are parallel. $\overline {XY}$ divides transversals $\overline {AB}$ and $\overline {CD}$ into a $1:1$ ratio (This can be shown by similar triangles.). Thus, it must also divide transversal $\overline {AC}$ and transversal $\overline {CO}$ into a $1:1$ ratio. By symmetry, the same applies for $\overline {CE}$ and $\overline {EA}$ as well as $\overline {EM}$ and $\overline {AN}.$

 In $\triangle ACE,$ we see that $\frac{[MNO]}{[ACE]} = \frac{1}{4}$ and $\frac{[MPN]}{[ACE]} = \frac{1}{8}.$ Our desired area becomes 
 \[\left(\frac{1}{4}+3 \cdot \frac{1}{8}\right) \cdot \frac{(\sqrt{3})^2 \cdot \sqrt{3}}{4} = \boxed{\textbf{(C)}\ \frac {15}{32}\sqrt{3}}.\]




\textbf{Solution 2 (Area Subtraction)}

Instead of directly finding the desired hexagonal area, $\triangle XYZ$ can be found. It consists of three triangles and the desired hexagon. Given triangle rotational symmetry, the three triangles are congruent. See that $\triangle XYZ$ and $\triangle ACE$ are equilateral, so $m\angle PXN=60,$ so $m\angle AXP = \frac{180-60}{2}=60.$ As $\overline {AC}$ is a transversal running through $\overline {FC}$ (use your imagination) and $\overline {AB},$ we have $m\angle BAC=m\angle FCA = \frac{m\angle ACE}{2}=30.$ 

 Then, $\triangle APX$ is a $30$-$60$-$90$ triangle. By HL congruence, $\triangle APX \cong \triangle NPX.$ Note that $AX=\frac{1}{2}.$ Then, the area of $\triangle PXN$ is $\frac{\sqrt{3}}{32}.$ There are three such triangles for a total area of $\triangle XYZ$ is $\frac{3\sqrt{3}}{32}.$ Find the side of $\triangle XYZ$ to be $\frac{3}{2},$ so the area is $\frac{9\sqrt{3}}{16}.$ 

 The answer is 
 \[\frac{9\sqrt{3}}{16}-\frac{3\sqrt{3}}{32}=\boxed{\textbf{(C)}\ \frac {15}{32}\sqrt{3}}.\]~BJHHar

 



\textbf{Solution 3 (Area Subtraction)}

We will find its area by subtracting the three triangles congruent to triangle APM from $\triangle ACE$. By the rotational symmetry of the regular hexagon, the three corner triangles outside the intersection ($\triangle APM$, $\triangle CQN$, and $\triangle ERO$) are all congruent. 

 We first find the area of the large equilateral triangle $\triangle ACE$. Since the regular hexagon has a side length of $1$, the side length of $\triangle ACE$ is $\sqrt{3}.$ Using the area formula for an equilateral triangle, we have:
 \[[\text{ACE}] = \frac{(\sqrt{3})^2\sqrt{3}}{4} = \frac{3\sqrt{3}}{4}.\]

 To find the area of the corner triangle $\triangle APM$, we can compare it to the smaller adjacent triangle $\triangle ZMR.$ Since $\triangle APM$ and $\triangle ZMR$ share the same vertical angles at vertex $M$ and are bounded by the grid lines of the hexagon, analyzing their bases and heights reveals that the area of $\triangle APM$ is exactly $3$ times the area of $\triangle ZMR.$ Thus, the area ratio between $\triangle APM$ and $\triangle ZMR$ is $3:1.$ 

 The area of the smaller triangle $\triangle ZMR$ can be calculated from the grid as:
 \[[\text{ZMR}] = \frac{1}{4} \cdot \frac{\sqrt{3}}{4} \cdot \frac{1}{2} = \frac{\sqrt{3}}{32}.\]

 Multiplying this by $3$ gives the true area of one corner triangle $\triangle APM$:
 \[[\text{APM}] = 3 \cdot \frac{\sqrt{3}}{32} = \frac{3\sqrt{3}}{32}.\]

 Since there are three such identical corner triangles surrounding the inner hexagon, we subtract their total area from the area of $\triangle ACE$:

 The answer is 
 \[\frac{3\sqrt{3}}{4} - 3\left(\frac{3\sqrt{3}}{32}\right) = \frac{24\sqrt{3}}{32} - \frac{9\sqrt{3}}{32} = \boxed{\textbf{(C)}\ \frac {15}{32}\sqrt{3}}.\]




\textbf{Solution 4 (Area Subtraction)}


                    \begin{figure}[H]
                        \centering
                        \includegraphics[width=0.5\textwidth]{Solution/AMC_12_12B_2018_20_Solution_Figure_1.png}
                    \end{figure}
                    

 Now, if we look at the figure, we can see that the complement of the hexagon we are trying to find is composed of $3$ isosceles trapezoids (namely $AXZF,CYXB,$ and $EZYD$) and $3$ right triangles (namely $\triangle XPN,\triangle YQO,$ and $\triangle ZRM$).

 Finding the trapezoid's area, we know that one base of each trapezoid is just the side length of the hexagon, which is $1,$ and the other base is $\frac{3}{2}$ (it is halfway in between the side and the longest diagonal, which has length $2$) with a height of $\frac{\sqrt{3}}{4}$ (by using the Pythagorean Theorem and the fact that it is an isosceles trapezoid) to give each trapezoid having an area of $\frac{5\sqrt{3}}{16}$ for a total area of $\frac{15\sqrt{3}}{16}.$ (Alternatively, we could have calculated the area of hexagon $ABCDEF$ and subtracted the area of $\triangle XYZ,$ which, as we showed before, had a side length of $\frac{3}{2}$).

 Now, we need to find the area of each of the small triangles, which, if we look at the triangle that has a vertex on $X,$ is similar to the triangle with a base of $YC = \frac12.$ Using similar triangles, we calculate the base to be $\frac{1}{4}$ and the height to be $\frac{\sqrt{3}}{4}$ giving us an area of  $\frac{\sqrt{3}}{32}$ per triangle, and a total area of $\frac{3\sqrt{3}}{32}.$ Adding the two areas together, we get $\frac{15\sqrt{3}}{16} + \frac{3\sqrt{3}}{32} = \frac{33\sqrt{3}}{32}.$ Finding the total area, we get $6 \cdot 1^2 \cdot \frac{\sqrt{3}}{4}=\frac{3\sqrt{3}}{2}.$ Taking the complement, we get $\frac{3\sqrt{3}}{2} - \frac{33\sqrt{3}}{32} = \boxed{\textbf{(C)}\ \frac {15}{32}\sqrt{3}}.$




\textbf{Solution 5 (Area Subtraction)}


                    \begin{figure}[H]
                        \centering
                        \includegraphics[width=0.5\textwidth]{Solution/AMC_12_12B_2018_20_Solution_Figure_3.png}
                    \end{figure}
                    

 We could also subtract $\triangle APM,\triangle CQN,$ and $\triangle ERO$ from $\triangle ACE.$

 Since $\angle BAF = 120^{\circ}$ and $\angle BAC = \angle FAE = 30^{\circ},$ we have $\angle CAE = \angle BAF-\angle BAC-\angle FAE=60^{\circ}.$

 Since $AX=BX$ and $FZ=EZ,$ we have $AF \parallel XZ,$ from which $\angle AMX= \angle FAM = 30^{\circ}.$

 We can show that $\triangle APM$ is $30$-$60$-$90$ using a similar method, $\triangle CQN$ and $\triangle ERO$ are also $30$-$60$-$90.$

 Since $AC=CE=AE=\sqrt{3},$ we have $[ACE]=AC^2 \cdot \frac{\sqrt{3}}{4}=3 \cdot \frac{\sqrt{3}}{4} = \frac{3 \sqrt{3}}{4}.$

 Since $AX= \frac{1}{2}$ and $AP = AX \cdot \frac{\sqrt{3}}{2}= \frac{1}{2} \cdot \frac{\sqrt{3}}{2} = \frac{\sqrt{3}}{4},$ we have $PM = AP \cdot \sqrt{3} = \frac{\sqrt{3}}{4} \cdot \sqrt{3} = \frac{3}{4}.$

 Note that 
 \[[APM]=[CQN]=[ERO]=\frac{1}{2} \cdot AP \cdot PM = \frac{1}{2} \cdot \frac{\sqrt{3}}{4} \cdot \frac{3}{4} = \frac{3 \sqrt{3}}{32}.\]Therefore, we get
 \[[PNQORM]=[ACE]-[APM]-[CQN]-[ERO]=\frac{3 \sqrt{3}}{4} - 3 \cdot \frac{3 \sqrt{3}}{32} = \frac{24 \sqrt{3}}{32} - \frac{9 \sqrt{3}}{32} = \boxed{\textbf{(C)}\ \frac {15}{32}\sqrt{3}}.\]




\textbf{Solution 6 (Partition the Hexagon)}

We partition hexagon $ABCDEF$ into $48$ congruent $30^\circ\text{-}60^\circ\text{-}90^\circ$ triangles, as shown below:
                    \begin{figure}[H]
                        \centering
                        \includegraphics[width=0.5\textwidth]{Solution/AMC_12_12B_2018_20_Solution_Figure_4.png}
                    \end{figure}
                    Let the brackets denote areas. Note that the desired region contains $15$ of the $48$ small triangles, so the answer is 
 \[\frac{15}{48}[ABCDEF]=\frac{15}{48}\cdot\frac{3\sqrt3}{2}=\boxed{\textbf{(C)}\ \frac {15}{32}\sqrt{3}}.\]~AlexLikeMath ~MRENTHUSIASM

 



\textbf{Solution 7 (Partition the Hexagon)}

Dividing $\triangle MNO$ into two right triangles congruent to $\triangle PMN,$ we see that $[MPNQOR]=\dfrac{5}{8}[ACE].$ Because $[ACE] = \dfrac{1}{2}[ABCDEF],$ we have $[MPNQOR]=\dfrac{5}{16}[ABCDEF].$ From here, you should be able to tell that the answer will have a factor of $5,$ and $\boxed{\textbf{(C)}\ \frac {15}{32}\sqrt{3}}$ is the only answer that has a factor of $5.$ However, if you want to actually calculate the area, you would calculate $[ABCDEF]$ to be $6 \cdot \dfrac{\sqrt{3}}{2 \cdot 2} = \dfrac{3\sqrt{3}}{2},$ so $[MPNQOR] = \dfrac{5}{16} \cdot \dfrac{3\sqrt{3}}{2} = \frac {15}{32}\sqrt{3}.$




\textbf{Solution 8 (Trigonometry)}

Notice, the area of the convex hexagon formed through the intersection of the $2$ triangles can be found by finding the area of the triangle formed by the midpoints of the sides and subtracting the smaller triangles that are formed by the region inside this triangle but outside the other triangle. First, let's find the area of the triangle formed by the midpoint of the sides. Notice, this is an equilateral triangle, thus all we need is to find the length of its side. To do this, we look at the isosceles trapezoid outside this triangle but inside the outer hexagon. Since the interior angle of a regular hexagon is $120^\circ$ and the trapezoid is isosceles, we know that the angle opposite is $60^\circ,$ and thus the side length of this triangle is $1+2\left(\frac{1}{2}\cos 60^\circ\right)=1+\frac{1}{2}=\frac{3}{2}.$ So the area of this triangle is $\frac{\sqrt{3}}{4}s^2=\frac{9\sqrt{3}}{16}.$

 Now let's find the area of the smaller triangles. Notice, $\triangle ACE$ cuts off smaller isosceles triangles from the outer hexagon. The base of these isosceles triangles is perpendicular to the base of the isosceles trapezoid mentioned before, thus we can use trigonometric ratios to find the base and height of these smaller triangles, which are all congruent due to the rotational symmetry of a regular hexagon. The area is then $\frac{1}{2}\left(\frac{1}{2}\cos 60^\circ\right)\left(\frac{1}{2}\sin 60^\circ\right)=\frac{\sqrt{3}}{32}$ and the sum of the areas is $3\cdot \frac{\sqrt{3}}{32}=\frac{3\sqrt{3}}{32}.$

 Therefore, the area of the convex hexagon is $\frac{9\sqrt{3}}{16}-\frac{3\sqrt{3}}{32}=\frac{18\sqrt{3}}{32}-\frac{3\sqrt{3}}{32}=\boxed{\textbf{(C)}\ \frac {15}{32}\sqrt{3}}.$




\textbf{Solution 9 (Linear Transformation)}

If we try to coordinate bash this problem, it's going to look very ugly with a lot of radicals. However, we can alter and skew the diagram in such a way that all ratios of lengths and areas stay the same while making it a lot easier to work with. Then, we can find the ratio of the area of the wanted region to the area of $ABCDEF$ then apply it to the old diagram.
                    \begin{figure}[H]
                        \centering
                        \includegraphics[width=0.5\textwidth]{Solution/AMC_12_12B_2018_20_Solution_Figure_5.png}
                    \end{figure}
                    

                    \begin{figure}[H]
                        \centering
                        \includegraphics[width=0.5\textwidth]{Solution/AMC_12_12B_2018_20_Solution_Figure_6.png}
                    \end{figure}
                    

 The isosceles right triangle with a leg length of $3$ in the new diagram is $\triangle XYZ$ in the old diagram. We see that if we want to take the area of the new hexagon, we must subtract $\frac{3}{4}$ from the area of $\triangle XYZ$ (the red triangles), giving us $\frac{15}{4}.$ However, we need to take the ratio of this area to the area of $ABCDEF,$ which is $\frac{\frac{15}{4}}{12}=\frac{5}{16}.$ Now we know that our answer is $\frac{5}{16} \cdot \frac{3\sqrt{3}}{2}=\boxed{\textbf{(C)}\ \frac {15}{32}\sqrt{3}}.$




\textbf{Solution 10 (Midsegment Theorem)}


                    \begin{figure}[H]
                        \centering
                        \includegraphics[width=0.5\textwidth]{Solution/AMC_12_12B_2018_20_Solution_Figure_7.png}
                    \end{figure}
                    (~MRENTHUSIASM's diagram + modifications)

 \textbf{Lemma 1: $MP$ is midsegment of $\triangle ANE$.} As $ZX$ is a midsegment of trapezoid $FEBA$, we know that $FA||ZX||EB$. Furthermore, by subdividing trapezoid $FEBA$, $ZM,MP,PX$ are midsegments of $\triangle EFA, \triangle ANE, \triangle, ANB$ respectively.

 \textbf{Lemma 2: $[ACE] = \frac12[ABCDEF]$} Due to radial symmetry about the center, $G$, we only have to prove $[AGC] = [CBA]$. As $\triangle AGB$ and $\triangle CBG$ are equilateral triangles, $AC$ therefore divides rhombus $ABCG$ into 2 congruent halves. $\triangle AGC$ and $\triangle CBA$ thus each contain two such congruent triangles, which means that their areas are equal.

 \textbf{Lemma 3: $[AMP] = \frac1{16}[ABCDEF]$} From Lemma 1, $\triangle AMP \simeq \triangle AEN$ with a scale factor of $\frac12$, which means $[AMP] = \frac14[AEN]$. Furthermore, due to symmetry of $\triangle AEN$ and $\triangle CEN$ about $EN$, $[AEN] = \frac12[ACE]$. Therefore, $[AMP] =  \frac18[ACE]$, or (using lemma 2) $\frac1{16}[ABCDEF]$.

 \textbf{Lemma 4: $[ABCDEF] = \frac{3\sqrt3}2$} Because the hexagon is regular, $[ABCDEF] = 6[ABG] = \frac{6\sqrt3}4 = \frac{3\sqrt3}2$. (For more info on where $\frac{\sqrt3}4$ came from, use the search term "area of equilateral triangle")

 \textbf{Final step:} To find the area of the shaded part, we seek $[ACE] - 3[AMP]$. Using lemmas 2 and 3, the expression resolves to $\frac12[ABCDEF] - 3\frac1{16}[ABCDEF] = \frac5{16} \cdot \frac{3\sqrt3}2 = \boxed{\textbf{(C)}\ \frac{15}{32}\sqrt{3}}$.

 


	\item (\textit{AMC 12 12B 2018}) In $\triangle{ABC}$ with side lengths $AB = 13$, $AC = 12$, and $BC = 5$, let $O$ and $I$ denote the circumcenter and incenter, respectively. A circle with center $M$ is tangent to the legs $AC$ and $BC$ and to the circumcircle of $\triangle{ABC}$. What is the area of $\triangle{MOI}$?


 $\textbf{(A)}\ \frac52\qquad\textbf{(B)}\ \frac{11}{4}\qquad\textbf{(C)}\ 3\qquad\textbf{(D)}\ \frac{13}{4}\qquad\textbf{(E)}\ \frac72$

 \textbf{Answer: }E



\textbf{Solution 1 (Coordinate Geometry)}

In this solution, let the brackets denote areas.

 We place the diagram in the coordinate plane: Let $A=(12,0),B=(0,5),$ and $C=(0,0).$

 Since $\triangle ABC$ is a right triangle with $\angle ACB=90^\circ,$ its circumcenter is the midpoint of $\overline{AB},$ from which $O=\left(6,\frac52\right).$ Note that the circumradius of $\triangle ABC$ is $\frac{13}{2}.$ 

 Let $s$ denote the semiperimeter of $\triangle ABC.$ The inradius of $\triangle ABC$ is $\frac{[ABC]}{s}=\frac{30}{15}=2,$ from which $I=(2,2).$

 Since $\odot M$ is also tangent to both coordinate axes, its center is at $M=(a,a)$ and its radius is $a$ for some positive number $a.$ Let $P$ be the point of tangency of $\odot O$ and $\odot M.$ As $\overline{OP}$ and $\overline{MP}$ are both perpendicular to the common tangent line at $P,$ we conclude that $O,M,$ and $P$ are collinear. It follows that $OM=OP-MP,$ or 
 \[\sqrt{(a-6)^2+\left(a-\frac52\right)^2}=\frac{13}{2}-a.\]Solving this equation, we have $a=4,$ from which $M=(4,4).$

 Finally, we apply the Shoelace Theorem to $\triangle MOI:$ 
 \[[MOI]=\frac12\left|\left(4\cdot\frac52+6\cdot2+2\cdot4\right)-\left(4\cdot6+\frac52\cdot2+2\cdot4\right)\right|=\boxed{\textbf{(E)}\ \frac72}.\]<u>\textbf{Remark}</u>

 Alternatively, we can use $\overline{MI}$ as the base and the distance from $O$ to $\overleftrightarrow{MI}$ as the height for $\triangle MOI:$

            \begin{itemize}
        \item By the Distance Formula, we have $MI=2\sqrt2.$
            \end{itemize}
            
            \begin{itemize}
        \item The equation of $\overleftrightarrow{MI}$ is $x-y+0=0,$ so the distance from $O$ to $\overleftrightarrow{MI}$ is $h_O=\frac{\left|1\cdot6+(-1)\cdot\frac52+0\right|}{\sqrt{1^2+(-1)^2}}=\frac74\sqrt2.$
            \end{itemize}
            Therefore, we get 
 \[[MOI]=\frac12\cdot MI\cdot h_O=\frac72.\]~pieater314159 ~MRENTHUSIASM

 



\textbf{Solution 2}

Let points Q and R be the points of tangency between the incircle and lines $AB$ and $BC$.  Notice that $\angle RIB$ is half of $\angle ABC$. Let $m\angle ABC = \theta$. Using the half angle tangent formula and remembering that $\cos(\theta) = \frac{5}{13}$, we find that $\tan(\theta/2) = \sqrt{\frac{1-\cos\theta}{1+\cos\theta}} = \sqrt{\frac{4}{9}} = \frac{2}{3}$. Now we can find the length of $BC$ in terms of the radius of the incircle, which will be $r$.  Using the knowledge that $\triangle RIB$ is right with some trigonometry, we find that $BR$ is $\frac{3r}{2}$. We also find that $CR$ is $r$, and so we can create the equation $BR + CR = BC \to \frac{3r}{2} + r = 5 \to \frac{5r}{2}$. We conclude $r = 2$. We now accept the fact that quadrilateral $BQIR$ is a kite, so $BR = BQ = 3$.  We also know that $O$ lies on $BA$ and divides $BA$ in half.  Next we determine that $MQ = BM - BQ = \frac{13}{2} - 3 = \frac{7}{2}$. We also know that $IQ$ has a measure of $2$ since it's the radius of circle $I$ 

 We know very little about the placement of $M$, so we work on that.  First, we can conclude that $\triangle IOQ$ is part of $\triangle MOI$. We guess that $M$ has to be higher up than $O$. We find the area of $\triangle IOQ$ to be $\frac{7}{2}$, which is the highest answer choice. I don't think that $M$ is below $\overline{AB}$, but if it were to go past $\overline{AB}$ then the answer would be GREATER than $\frac{7}{2}$. This means $\overline{QO}$ is on $\overline{AB}$ . Therefore $\triangle IOQ=\triangle MOI$. The answer must be $\boxed{\textbf{(E)}\ \frac72}$ .

 



\textbf{Solution 3}

The circle with center $M$ is the $C$-mixtilinear incircle $\omega$ of $\triangle ABC$. Let $T$ be the intersection between $\omega$ and the circumcircle $\Omega$. Then, there is a homothety centered at $T$ sending $\omega$ to $\Omega$. As such, the tangent $\overline{AC}$ gets sent to a parallel tangent to $\Omega$, which thus must be the arc midpoint $M_B$ of arc ${CA}$. Thus, by inscribed angle theorem $CM_C$ and $BM_B$ intersect at $I$, so $EF$ passes through $I$ where $E$, $F$ are the tangency points of $\omega$ with $\overline{BC}$ and $\overline{AC}$ by Pascal's Theorem. Thus, we see since $\angle BCA = 90 ^{\circ}$, $r_{\omega} = 2 r_{\text{incircle}} = 4$. Set up a coordinate plane and apply Shoelace to obtain $[MOI] = \boxed{ \textbf{(E)} \frac{7}{2}}$.

 


	\item (\textit{AMC 12 12B 2018}) Consider polynomials $P(x)$ of degree at most $3$, each of whose coefficients is an element of $\{0, 1, 2, 3, 4, 5, 6, 7, 8, 9\}$. How many such polynomials satisfy $P(-1) = -9$?


 $\textbf{(A) } 110 \qquad \textbf{(B) } 143 \qquad \textbf{(C) } 165 \qquad \textbf{(D) } 220 \qquad \textbf{(E) } 286$

 \textbf{Answer: }D



\textbf{Solution 1 (Stars and Bars)}

Suppose that $P(x)=ax^3+bx^2+cx+d.$ This problem is equivalent to counting the ordered quadruples $(a,b,c,d),$ where all of $a,b,c,$ and $d$ are integers from $0$ through $9$ such that 
 \[P(-1)=-a+b-c+d=-9.\] Let $a'=9-a$ and $c'=9-c.$ Note that both of $a'$ and $c'$ are integers from $0$ through $9.$ Moreover, the ordered quadruples $(a,b,c,d)$ and the ordered quadruples $(a',b,c',d)$ have one-to-one correspondence.

 We rewrite the given equation as $(9-a)+b+(9-c)+d=9,$ or 
 \[a'+b+c'+d=9.\] By the stars and bars argument, there are $\binom{9+4-1}{4-1}=\boxed{\textbf{(D) } 220}$ ordered quadruples $(a',b,c',d).$




\textbf{Solution 2 (Casework)}

Suppose that $P(x)=ax^3+bx^2+cx+d.$ This problem is equivalent to counting the ordered quadruples $(a,b,c,d),$ where all of $a,b,c,$ and $d$ are integers from $0$ through $9$ such that $P(-1)=-a+b-c+d=-9,$ which rearranges to 
 \[b+d+9=a+c.\]Note that $b+d+9$ is an integer from $9$ through $27,$ and $a+c$ is an integer from $0$ through $18.$ Therefore, both of $b+d+9$ and $a+c$ are integers from $9$ through $18.$ We construct the following table:
 \[\begin{array}{c|c|c|c||c} & & & & \\ [-2.5ex] \boldsymbol{b+d} & \boldsymbol{\#}\textbf{ of Ordered Pairs }\boldsymbol{(b,d)} & \boldsymbol{a+c} & \boldsymbol{\#}\textbf{ of Ordered Pairs }\boldsymbol{(a,c)} & \boldsymbol{\#}\textbf{ of Ordered Quadruples }\boldsymbol{(a,b,c,d)} \\ [0.5ex] \hline & & & & \\ [-2ex] 0 & 1 & 9 & 10 & 1\cdot10=10 \\ 1 & 2 & 10 & 9 & \phantom{0}2\cdot9=18 \\ 2 & 3 & 11 & 8 & \phantom{0}3\cdot8=24 \\ 3 & 4 & 12 & 7 & \phantom{0}4\cdot7=28 \\ 4 & 5 & 13 & 6 & \phantom{0}5\cdot6=30 \\ 5 & 6 & 14 & 5 & \phantom{0}6\cdot5=30 \\ 6 & 7 & 15 & 4 & \phantom{0}7\cdot4=28 \\ 7 & 8 & 16 & 3 & \phantom{0}8\cdot3=24 \\ 8 & 9 & 17 & 2 & \phantom{0}9\cdot2=18 \\ 9 & 10 & 18 & 1 & 10\cdot1=10 \end{array}\]We sum up the counts in the last column to get the answer $2\cdot(10+18+24+28+30)=\boxed{\textbf{(D) } 220}.$



	\item (\textit{AMC 12 12B 2018}) Ajay is standing at point $A$ near Pontianak, Indonesia, $0^\circ$ latitude and $110^\circ \text{ E}$ longitude. Billy is standing at point $B$ near Big Baldy Mountain, Idaho, USA, $45^\circ \text{ N}$ latitude and $115^\circ \text{ W}$ longitude. Assume that Earth is a perfect sphere with center $C.$ What is the degree measure of $\angle ACB?$
$\textbf{(A) }105 \qquad \textbf{(B) }112\frac{1}{2} \qquad \textbf{(C) }120 \qquad \textbf{(D) }135 \qquad \textbf{(E) }150 \qquad$

 \textbf{Answer: }C



\textbf{Solution 1 (Tetrahedron)}

This solution refers to the \textbf{Diagram} section.

 Let $D$ be the orthogonal projection of $B$ onto the equator. Note that $\angle BDA = \angle BDC = 90^\circ$ and $\angle BCD = 45^\circ.$ Recall that $115^\circ \text{ W}$ longitude is the same as $245^\circ \text{ E}$ longitude, so $\angle ACD=135^\circ.$

 We obtain the following diagram:
                    \begin{figure}[H]
                        \centering
                        \includegraphics[width=0.5\textwidth]{Solution/AMC_12_12B_2018_23_Solution_Figure_1.png}
                    \end{figure}
                    Without loss of generality, let $AC=BC=1.$ For tetrahedron $ABCD:$

            \begin{enumerate} 
        \item Since $\triangle BCD$ is an isosceles right triangle, we have $BD=CD=\frac{\sqrt2}{2}.$

        \item In $\triangle ACD,$ we apply the Law of Cosines to get $AD=\sqrt{AC^2+CD^2-2\cdot AC\cdot CD\cdot\cos\angle ACD}=\frac{\sqrt{10}}{2}.$

        \item In right $\triangle ABD,$ we apply the Pythagorean Theorem to get $AB=\sqrt{AD^2+BD^2}=\sqrt{3}.$

        \item In $\triangle ABC,$ we apply the Law of Cosines to get $\cos\angle ACB=\frac{AC^2+BC^2-AB^2}{2\cdot AC\cdot BC}=-\frac12,$ so $\angle ACB=\boxed{\textbf{(C) }120}$ degrees.

            \end{enumerate}
            



\textbf{Solution 2 (Cartesian Coordinates and Vectors)}

This solution refers to the \textbf{Diagram} section.

 Let $D$ be the orthogonal projection of $B$ onto the equator. Note that $\angle BDA = \angle BDC = 90^\circ$ and $\angle BCD = 45^\circ.$ Recall that $115^\circ \text{ W}$ longitude is the same as $245^\circ \text{ E}$ longitude, so $\angle ACD=135^\circ.$

 Without loss of generality, let $AC=BC=1.$ As shown below, we place Earth in the $xyz$-plane with $C=(0,0,0)$ such that the positive $x$-axis runs through $A,$ the positive $y$-axis runs through $0^\circ$ latitude and $160^\circ \text{ W}$ longitude, and the positive $z$-axis runs through the North Pole.
                    \begin{figure}[H]
                        \centering
                        \includegraphics[width=0.5\textwidth]{Solution/AMC_12_12B_2018_23_Solution_Figure_2.png}
                    \end{figure}
                    It follows that $A=(1,0,0)$ and $D=(-t,t,0)$ for some positive number $t.$ Since $\triangle BCD$ is an isosceles right triangle, we have $B=\left(-t,t,\sqrt{2}t\right).$ By the Distance Formula, we get $(-t)^2+t^2+\left(\sqrt{2}t\right)^2=1,$ from which $t=\frac12.$

 As $\vec{A} = \begin{pmatrix}1 \\ 0 \\ 0 \end{pmatrix}$ and $\vec{B} = \begin{pmatrix}-1/2 \\ 1/2 \\ \sqrt2/2 \end{pmatrix},$ we obtain 
 \[\cos\angle ACB=\frac{\vec{A}\bullet\vec{B}}{\left\lVert\vec{A}\right\rVert\left\lVert\vec{B}\right\rVert}=-\frac12\] by the dot product, so $\angle ACB=\boxed{\textbf{(C) }120}$ degrees.

 



\textbf{Solution 3 (Spherical Coordinates and Vectors)}

This solution refers to the diagram in Solution 2.

 In spherical coordinates $(\rho,\theta,\phi),$ note that $\rho,\theta,$ and $\phi$ represent the radial distance, the polar angle, and the azimuthal angle, respectively.

 Without loss of generality, let $AC=BC=1.$ As shown in Solution 2, we place Earth in the $xyz$-plane with origin $C$ such that the positive $x$-axis runs through $A,$ the positive $y$-axis runs through $0^\circ$ latitude and $160^\circ \text{ W}$ longitude, and the positive $z$-axis runs through the North Pole.

 In spherical coordinates, we have $A=(1,90^\circ,0^\circ)$ and $B=(1,45^\circ,135^\circ).$ Now, we express $\vec{A}$ and $\vec{B}$ in Cartesian coordinates: 
 \[\vec{A} = \begin{pmatrix}\sin90^\circ \cos0^\circ \\ \sin90^\circ \sin0^\circ \\ \cos90^\circ \end{pmatrix} = \begin{pmatrix}1 \\ 0 \\ 0 \end{pmatrix} \text{ and } \vec{B} = \begin{pmatrix}\sin45^\circ \cos135^\circ \\ \sin45^\circ \sin135^\circ \\ \cos45^\circ \end{pmatrix} = \begin{pmatrix}-1/2 \\ 1/2 \\ \sqrt2/2 \end{pmatrix}.\]We continue with the last paragraph of Solution 2 to get the answer $\angle ACB=\boxed{\textbf{(C) }120}$ degrees.

 



\textbf{Solution 4 (It's Not That Deep)}

This solution refers to the diagram in Solution 2. 

 Let the radius of the sphere be $r$. Form $\triangle{ABC}$ and notice that it is isosceles, with $AC=BC$. Draw the circle with radius equal to the sphere (great circle) with $0^{\circ}$ latitude. 

 Drop a perpendicular from $B$ onto this circle. Let the foot of this perpendicular be called $D$. Since $\angle{DCA} = 135^{\circ}$, by Law of Cosines on $\triangle{ACD}$, $(AD)^2 = \frac{5}{2}r^2$. Since $BD$ is a perpendicular, we can form right triangle $\triangle{BDA}$. By Pythag, $AB=\sqrt{\frac{1}{2}r^2 + \frac{5}{2}r^2} = r\sqrt{3}$. Since $\triangle{ABC}$ has side lengths $r, r,$ and $r\sqrt{3}$, $\angle{BCA} = \boxed{\textbf{(C) }120^{\circ}}$.

 -skibbysiggy

 


	\item (\textit{AMC 12 12B 2018}) Let $\lfloor x \rfloor$ denote the greatest integer less than or equal to $x$. How many real numbers $x$ satisfy the equation $x^2 + 10,000\lfloor x \rfloor = 10,000x$?


 $\textbf{(A) } 197 \qquad \textbf{(B) } 198 \qquad \textbf{(C) } 199 \qquad \textbf{(D) } 200 \qquad \textbf{(E) } 201$

 \textbf{Answer: }C



\textbf{Solution 1}

This rewrites itself to $x^2=10,000\{x\}$ where $\lfloor x \rfloor + \{x\} = x$.

 Graphing $y=10,000\{x\}$ and $y=x^2$ we see that the former is a set of line segments with slope $10,000$ from $0$ to $1$ with a hole at $x=1$, then $1$ to $2$ with a hole at $x=2$ etc.Here is a graph of $y=x^2$ and $y=16\{x\}$ for visualization.

 
                    \begin{figure}[H]
                        \centering
                        \includegraphics[width=0.5\textwidth]{Solution/AMC_12_12B_2018_24_Solution_Figure_1.png}
                    \end{figure}
                    

 Now notice that when $x=\pm 100$ the graph has a hole at $(\pm 100,10,000)$ which the equation $y=x^2$ passes through and then continues upwards. Thus our set of possible solutions is bounded by $(-100,100)$. We can see that $y=x^2$ intersects each of the lines once and there are $99-(-99)+1=199$ lines for an answer of $\boxed{\text{(C)}~199}$.

 



\textbf{Solution 2}

Same as the first solution, $x^2=10,000\{x\}$.

 
We can write $x$ as $\lfloor x \rfloor+\{x\}$. Expanding everything, we get a quadratic in $\{x\}$ in terms of $\lfloor x \rfloor$:
 \[\{x\}^2+ (2\lfloor x \rfloor -10,000)\{x\} + \lfloor x \rfloor ^2 = 0\]

 
We use the quadratic formula to solve for $\{x\}$ :
 \[\{x\} = \frac {-2\lfloor x \rfloor + 10,000 \pm \sqrt{ ( 2\lfloor x \rfloor - 10,000 )^2 - 4\lfloor x \rfloor^2  }}{2} = \frac {-2\lfloor x \rfloor + 10,000 \pm \sqrt{ 4\lfloor x \rfloor^2 -40,000 \lfloor x \rfloor + 10,000^2- 4\lfloor x \rfloor^2  }}{2}\]

 
Since $0 \leq \{x\} < 1$, we get an inequality which we can then solve. After simplifying a lot, we get that $\lfloor x \rfloor^2 + 2\lfloor x \rfloor - 9999 < 0$.

 
Solving over the integers, $-101 < \lfloor x \rfloor < 99$, and since $\lfloor x \rfloor$ is an integer, there are $\boxed{\text{(C)}~199}$ solutions. Each value of $\lfloor x \rfloor$ should correspond to one value of $x$, so we are done.

 



\textbf{Solution 3 (Fakesolve)}

Let $x = a+k$ where $a$ is the integer part of $x$ and $k$ is the fractional part of $x$.We can then rewrite the problem below:

 $(a+k)^2 + 10000a = 10000(a+k)$

 From here, we get

 $(a+k)^2 + 10000a = 10000a + 10000k$

 Solving for $a+k = x$
$(a+k)^2 = 10000k$
$x = a+k = \pm100\sqrt{k}$

 Because $0 \leq k < 1$, we know that $a+k$ cannot be less than or equal to $-100$ nor greater than or equal to $100$. Therefore:

 $-99 \leq x \leq 99$

 There are $199$ elements in this range, so the answer is $\boxed{\textbf{(C)} \text{ 199}}$.

 
Note (not by author): this solution seems to be invalid at first, because one can not determine whether $x$ is an integer or not. However, it actually works because although $x$ itself might not be an integer, it is very close to one, so there are 199 potential $x$.

 
Another Note (not by author of previous note): we can actually determine that $x$=0 is the only possible integer value of $x$ is we set $x$=$\lfloor x \rfloor$ we end up with $x$=0   ~YJC64002776

 
Yet Another Note EXCEPT THIS ONE'S VERY IMPORTANT, AS IT DISPROVES THE PROOF (also not by the author of any previous notes or of the original solution): this solution claims that $-99 \leq x \leq 99$, which is not entirely true. When $x = 5000-1000\sqrt{26}$, the equation originally stated in the question holds true. $5000-1000\sqrt{26}$ is roughly $-99.0195$ (rounded to the nearest ten thousandth), which is not in the $-99 \leq x \leq 99$ interval. It is probably by mere coincidence that this solution gives the correct answer, as the solution is not rigorous. Additionally, $x$ is not necessarily integer, so it is not claimable that there are 199 elements in the range of $-99 \leq x \leq 99$.

 
One More Note (not by any of the previous authors): The original guy who wrote the solution just wrote the wrong range lol this is a valid solution (at least I'm pretty sure). Given $(a+k)^2 = 10000k$, dividing each side by 10000 gives $\frac{(a+k)^2}{10000} = k$. Our original definition of $k$ states that $0 \leq k < 1$, so $0 \leq \frac{(a+k)^2}{10000} < 1$. This is equivalent to the range $-100 \leq x < 100$ NOT $-99 \leq x < 99$. If you observe the functions behavior, you realize that between every interval of length $1$ there is one solution, except for the intervals that include $0$, where only $0$ is a solution. There are therefore $99+99+1$ solutions in that range which gives the answer. Also $x$ literally doesn't have to be an integer the original constraint was that $x$ is a real number.

 



\textbf{Solution 4}

Notice the given equation is equivalent to $(\lfloor x \rfloor+\{x\})^2=10,000\{x\}$

 Now we know that $\{x\} < 1$ so plugging in $1$ for $\{x\}$ we can find the upper and lower bounds for the values.

 $(\lfloor x \rfloor +1)^2 = 10,000(1)$
$(\lfloor x \rfloor +1) = \pm 100$
$\lfloor x \rfloor = 99, -101$

 And just like $\textbf{Solution 2}$, we see that $-101 < \lfloor x \rfloor < 99$, and since $\lfloor x \rfloor$ is an integer, there are $\boxed{\text{(C)}~199}$ solutions. Each value of $\lfloor x \rfloor$ should correspond to one value of $x$, so we are done.

 
            \begin{itemize}
        \item $\textbf{Note}$ (not by author of any other notes beforehand): this solution is probably the best and most efficient of all. It solves the problems stated in the fakesolves./
            \end{itemize}
            



\textbf{Solution 5}

Firstly, if $x$ is an integer, then $10,000\lfloor x \rfloor=10,000x$, so $x$ must be $0$.

 If $0<x<1$, then we know the following:

 $0<x^2<1$
$10,000\lfloor x \rfloor =0$
$0<10,000x<10,000$

 Therefore, $0<x^2+10,000\lfloor x \rfloor <1$, which overlaps with $0<10,000x<10,000$. This means that there is at least one real solution between $0$ and $1$. Since $x^2+10,000\lfloor x \rfloor$ increases quadratically and $10,000x$ increases linearly, there is only one solution for this case. 

 Similarly, if $1<x<2$, then we know the following:

 $1<x^2<4$
$10,000\lfloor x \rfloor =10,000$
$<10,000<10,000x<20,000$

 By following similar logic, we can find that there is one solution between $1$ and $2$. 

 We can also follow the same process to find that there are negative solutions for $x$ as well.

 There are not an infinite amount of solutions, so at one point there will be no solutions when $n<x<n+1$ for some integer $n$. For there to be no solutions in a given range means that the range of $10,000\lfloor x \rfloor + x^2$ does not intersect the range of $10,000x$. $x^2$ will always be positive, and $10,000\lfloor x \rfloor$ is less than $10,000$ less than $10,000x$, so when $x^2 >= 10,000$, the equation will have no solutions. This means that there are $99$ positive solutions, $99$ negative solutions, and $0$ for a total of $\boxed{\text{(C)}~199}$ solutions.

 



\textbf{Solution 6 (General Equation)}

General solution to this type of equation $f(x, \lfloor x \rfloor) = 0$:

 \begin{verbatim}
1. solve $f(x, \lfloor x \rfloor) = 0$ for $x$ to get $x = g(\lfloor x \rfloor )$2. apply $\lfloor x \rfloor \le x < \lfloor x \rfloor+1$, solve $\lfloor x \rfloor \le g(\lfloor x \rfloor) < \lfloor x \rfloor+1$ to get the domain of $\lfloor x \rfloor$3. get $\lfloor x \rfloor$ from the domain of $\lfloor x \rfloor$ because $\lfloor x \rfloor$ is integer, then get $x$ from $\lfloor x \rfloor$ by $x = g( \lfloor x \rfloor)$Note: function $\lfloor x \rfloor$ maps $x$ to its floor. By solving $f(x, \lfloor x \rfloor) = 0$, we get function $x = g( \lfloor x \rfloor)$, mapping $x$'s floor to $x$
\end{verbatim}

$x^2 - 10000x + 10000 \lfloor x \rfloor =0$
$x=5000 \pm 100 \sqrt{2500- \lfloor x \rfloor}$, $\lfloor x \rfloor \le 2500$
$\lfloor x \rfloor \le x < \lfloor x \rfloor + 1$

 If $x= 5000 + 100 \sqrt{2500 - \lfloor x \rfloor}$, $x \ge 5000$, it contradicts $x < \lfloor x \rfloor + 1 \le 2501$

 So $x= 5000 - 100 \sqrt{2500 - \lfloor x \rfloor}$

 Let $k = \lfloor x \rfloor$ , $x= 5000 - 100 \sqrt{2500 - k}$
$k \le 5000 - 100 \sqrt{2500 - k} < k + 1$
$0 \le 5000 - k - 100 \sqrt{2500 - k} < 1$
$0 \le 2500 - k - 100 \sqrt{2500 - k} + 2500 < 1$
$0 \le (\sqrt{2500 - k} - 50)^2 < 1$
$-1 < \sqrt{2500 - k} - 50 < 1$
$49 < \sqrt{2500 - k} < 51$
$-101 < k < 99$

 So the number of $k$'s values is $99-(-101)-1=199$. Because $x=5000-100\sqrt{2500-k}$, for each value of $k$, there is a value for $x$. The answer is $\boxed{\textbf{(C)} 199}$




\textbf{Solution 7 (Fakesolve)}

Subtracting $10000\lfloor x\rfloor$ from both sides gives $x^2=10000(x-\lfloor x\rfloor)=10000\{x\}$. Dividing both sides by $10000$ gives $\left(\frac{x}{100}\right)^2=\{x\}<1$.  $\left(\frac{x}{100}\right)^2<1$ when $-100<x<100$ so the answer is $\boxed{199}$.

 



\textbf{Solution 8 (Also Gives General Formula For Values of x)}

The question wants to know how many values of $x$ satisfy the equation $x^2 + 10,000\lfloor x \rfloor = 10,000x$. This equation can be simplified as follows

 
 \[x^2 + 10,000\lfloor x \rfloor = 10,000x\]

 \[x^2 = 10,000(x-\lfloor x \rfloor)\]

 Notice that $x-\lfloor x \rfloor$ must be greater than or equal to zero, but less than one. Because of that, the right-hand side of the equation must be less than $10000$. Therefore, $-100<x<100$.

 $x$ can be expressed as $a+b$, where $a$ is an integer and $b$ is a real number such that $0 \leq b < 1$. Notice that in order to satisfy the conditions set down for $x$, $-100 \leq a \leq 99$.

 Substituting $a$ and $b$ into $x^2 = 10000(x-\lfloor x \rfloor)$, we get

 
 \[(a+b)^2 = 10,000b\]

 \[a^2+2ab+b^2 = 10,000b\]

 Let's try to turn this into a quadratic equation where we're trying to solve for $b$. Simplifying, we find

 
 \[b^2+(2a-10,000)b+a^2 = 0\]

 Now the value of $b$ depends on the value of $a$. Our task is to figure out how many values of $a$ will give me valid values of $b$ (a.k.a the values of $a$ that give me $b$ such that $0 \leq b < 1$), as that will be our answer. We must also keep in mind that $-100 \leq a \leq 99$.

 Using the quadratic formula, we can find

 
 \[\frac{(-2a+10,000) \pm \sqrt{(2a-10,000)^2-4a^2}}{2}\]

 The expression above represents the values of $b$ in terms of $a$. Since it represents $b$, we want values of $a$ such that the equation is between zero and one, zero inclusive. We can then conceive the inequality below.

 
 \[0 \leq \frac{(-2a+10,000) \pm \sqrt{(2a-10,000)^2-4a^2}}{2} < 1\]

 Wow, that looks like a mess! How are we supposed to easily use such a complicated expression? Maybe it would help to simplify it further. First, let's notice a few things.

 
            \begin{itemize}
        \item I want the numerator to be between 0 and 2, 0 inclusive. Of the $\pm$ in the quadratic formula, only the "$-$" (minus) will work. Why? Because $-2a+10000$ will always be positive (if I plug in the largest possible value of $a$, $99$, it will give $9802$), and if I want $b$ to be real, which I do, then $\sqrt{(2a-10,000)^2-4a^2}$ should also be positive. $-2a+10000$ will already be much bigger than 2, and adding more on will only make it bigger. It will never be less than 2. It will never work.
            \end{itemize}
            
            \begin{itemize}
        \item I can simplify the numbers under the square root using Difference of Squares. I can turn $\sqrt{(2a-10,000)^2-4a^2}$ into $\sqrt{(2a-10,000)^2-(2a)^2}$. Applying Difference of Squares, it will eventually simplify out to be $200\sqrt{2500-a}$.
            \end{itemize}
            Cool! Now we've simplified the square root and figured out that it must be a minus sign, not a plus sign. Our new inequality looks like this:

 
 \[0 \leq \frac{(-2a+10,000) - 200\sqrt{2500-a}}{2} < 1\]

 Simplifying the unsimplified fraction further, we can get this:

 
 \[0 \leq -a + 5,000 - 100\sqrt{2500-a} < 1\]

 Yay! A simpler equation. As long as the value $a$ fits this inequality, the restrictions set on $b$ are satisfied. But wait. What about the restrictions on $a$?

 \textbf{By testing out $-100$ and $99$ (and finding that $-100$ gives a value very close to $1$, and that $99$ gives exactly $1$), we can conclude that any value of $-100 \leq$a$\leq 99$ will satisfy the inequality above as long as it does not make $-a + 5,000 - 100\sqrt{2500-a}$ equal to $1$. By testing $-100$ and $99$, you will also find that $a = 99$ will not satisfy the inequality.} <u><-- old, non-rigorous way. You can ignore it, but if you want to see a bogus way to prove it, feel free to read it.</u>

 I can solve $0 \leq -a + 5000 - 100\sqrt{2500-a} < 1$ for $a$ this way:

 $0 \leq -a + 5000 - 100\sqrt{2500-a} < 1$
$-5000 \leq -a - 100\sqrt{2500-a} < -4999$
$5000 \geq a + 100\sqrt{2500-a} < 4999$
$5000-a \geq 100\sqrt{2500-a} < 4999-a$
$25000000-10000a+a^2 \geq 25000000-10000a > 4999^2-9998a+a^2$
$a^2 \geq 0 > -9999+2a+a^2$
$0 \geq -a^2 > -9999+2a$
$0 \leq a^2 < 9999-2a$
$2a \leq a(a+2) < 9999$

 Following the left part of the inequality,

 $2a \leq a^2 + 2a$
$0 \leq a^2$
$a$ can be anything $(-\infty, \infty)$.

 Following the right part of the inequality,

 $a(a+2) < 9999$
$a < 99$ and $a > -101$
$a$ can be anything $(-101, 99)$.

 Combining both restrictions and given previously that $-100 \leq a \leq 99$, we can conclude that the inequality describing all possible values of $a$ is $-100 \leq a < 99$.

 Remember how these values of $a$ give valid values of $b$ that satisfy the desired inequality $0 \leq b < 1$ and how adding up $a$ and $b$ gives $x$. In summary, remember that the number of values of $a$ is equivalent to the number of values of $x$.

 $a$ must be integer, so the number of values of $a$ that satisfy $-100 \leq a < 99$ is $\boxed{199}$.

 The answer is $\boxed{\textbf{(C)} 199}$.

 
Note From Author of Solution: you may have noticed that I missed a slight bit by not proving that the thing under the square root in the quadratic formula I made was positive. Don't worry, it wasn't a mistake. It's actually pretty easy to prove. I'll leave it to you to figure out how to do it!

 



\textbf{Solution 9}

Note that the problem rearranges itself to $x^2=10,000\{x\}$ where $\lfloor x \rfloor + \{x\} = x$.

 
To get a general scope of solutions, by noting that $0 \leq \{x\} < 1$, we have that $0 \leq x^2 < 10,000$. Therefore, $-100 < x < 100$.

 
From here, it is easy to observe that there is one solution per each integer value of $-99 \leq \lfloor x \rfloor \leq 99$. This can be reasoned by bounding $\{x\}$ for each $\lfloor x \rfloor$ (intuitively, this is because $\lfloor x \rfloor$ determines the fractional part: for each $\lfloor x \rfloor$, you will be forced to choose the value of $\{x\}$ that will satisfy the equation).

 
Ultimately, the answer is $99+99+1 = \boxed{\textbf{(C)} 199}$.

 


 


	\item (\textit{AMC 12 12B 2018}) Circles $\omega_1$, $\omega_2$, and $\omega_3$ each have radius $4$ and are placed in the plane so that each circle is externally tangent to the other two.  Points $P_1$, $P_2$, and $P_3$ lie on $\omega_1$, $\omega_2$, and $\omega_3$ respectively such that $P_1P_2=P_2P_3=P_3P_1$ and line $P_iP_{i+1}$ is tangent to $\omega_i$ for each $i=1,2,3$, where $P_4 = P_1$.  See the figure below.  The area of $\triangle P_1P_2P_3$ can be written in the form $\sqrt{a}+\sqrt{b}$ for positive integers $a$ and $b$.  What is $a+b$?


 
                    \begin{figure}[H]
                        \centering
                        \includegraphics[width=0.5\textwidth]{Problem/AMC_12_12B_2018_25_Problem_Figure_1.png}
                    \end{figure}
                    
$\textbf{(A) }546\qquad\textbf{(B) }548\qquad\textbf{(C) }550\qquad\textbf{(D) }552\qquad\textbf{(E) }554$

 \textbf{Answer: }D



\textbf{Solution 1}

Let $O_1$ and $O_2$ be the centers of $\omega_1$ and $\omega_2$ respectively and draw $O_1O_2$, $O_1P_1$, and $O_2P_2$. Note that $\angle{O_1P_1P_2}$ and $\angle{O_2P_2P_3}$ are both right. Furthermore, since $\triangle{P_1P_2P_3}$ is equilateral, $m\angle{P_1P_2P_3} = 60^\circ$ and $m\angle{O_2P_2P_1} = 30^\circ$. Mark $M$ as the base of the altitude from $O_2$ to $P_1P_2.$ Since $\triangle P_2O_2M$ is a 30-60-90 triangle, $O_2M = 2$ and $P_2M = 2\sqrt{3}$. Also, since $O_1O_2 = 8$ and $O_1P_1 = 4$, we can find $P_1M = \sqrt{8^2 - (4 + 2)^2} = 2\sqrt{7}$. Thus, $P_1P_2 = P_1M + P_2M = 2\sqrt{3} + 2\sqrt{7}$. This makes 
 \[\left[P_1P_2P_3\right] = \frac{\sqrt 3}4\cdot\left(2\sqrt 3 + 2\sqrt 7\right)^2 = 10\sqrt 3 + 6\sqrt 7 = \sqrt{300} + \sqrt{252}.\] So, our answer is $252 + 300 = \boxed{\textbf{D) }552}$.

 Note by diyarv:

 A way to see that $P_1M = 2{\sqrt 7}$ is through similar triangles. Call the intersection between $P_1M$ and $O_1O_2$ $S$. Since angles $\angle{O_1P_1S}$ and $\angle{O_2MS}$ are both right angles, and $\angle{P_1SO_1}$ and $\angle{MSO_2}$ are congruent, triangles $\triangle{P_1SO_1}$ and $\triangle{MSO_2}$ are similar by AA similarity. And, since $O_1P_1$ is $4$ and $O_2M$ is 2, the common ratio is $2$. Using the fact that $O_1S + SO_2 = 8$, we see that these lengths are $\frac{16}{3}$ and $\frac{8}{3}$ respectively. Using the Pythagorean theorem, we see that $P_1S$ and $SP_2$ are $\frac{2\sqrt{7}}{3}$ and $\frac{4\sqrt{7}}{3}$ respectively, giving us a sum of $2\sqrt{7}$.

 A better note on why $P_1M = \sqrt{8^2 - (4 + 2)^2}$:

 This is more easily seen by translating $O_1O_2$ so that $O_2$ is mapped to $M$ and the resulting $O_1’$ remains collinear with the original $O_1$ and $P_1$. This forms a right triangle, with leg $O_1’ P = 4+2 = 6$ and hypotenuse $O_1’O_2 = 8$. Therefore $P_1M = \sqrt{8^2 - (4 + 2)^2}$




\textbf{Solution 2}

Let $O_i$ be the center of circle $\omega_i$ for $i=1,2,3$, and let $K$ be the intersection of lines $O_1P_1$ and $O_2P_2$. Because $\angle P_1P_2P_3 = 60^\circ$, it follows that $\triangle P_2KP_1$ is a $30-60-90$ triangle. Let $x=P_1K$; then $P_2K = 2x$ and $P_1P_2 = x \sqrt 3$. The Law of Cosines in $\triangle O_1KO_2$ gives 
 \[8^2 = (x+4)^2 + (2x-4)^2 - 2(x+4)(2x-4)\cos 60^\circ,\]which simplifies to $3x^2 - 12x - 16 = 0$. The positive solution is $x = 2 + \frac23\sqrt{21}$. Then $P_1P_2 = x\sqrt 3 = 2\sqrt 3 + 2\sqrt 7$, and so the area of $\triangle P_1P_2P_3$ is 
 \[\frac{\sqrt 3}4\cdot\left(2\sqrt 3 + 2\sqrt 7\right)^2 = 10\sqrt 3 + 6\sqrt 7 = \sqrt{300} + \sqrt{252}.\]The requested sum is $300 + 252 = \boxed{552}$.

 



\textbf{Solution 3}

Let $O_i$ be the center of circle $\omega_i$ for $i=1,2,3$. Let $X$ be the centroid of $\triangle{O_1O_2O_3}$, which also happens to be the centroid of $\triangle{P_1P_2P_3}$. Because $m\angle{O_1P_1P_2} = 90^\circ$ and $m\angle{O_1P_1X} = 30^\circ$, $m\angle{O_1P_1X} = 60^\circ$. $O_2M$ is $2/3$ the height of  $\triangle{P_1P_2P_3}$, thus $O_2M$  is  $8*\sqrt{3}/3$.

 Applying cosine law on $\triangle{O_1P_1X}$, one finds that $P_1X = 2 + 2*\sqrt{21}/3$. Multiplying by $3/2$ to solve for the height of $\triangle{P_1P_2P_3}$, one gets $3 + \sqrt{21}$. Simply multiplying by $2/\sqrt{3}$ and then calculating the equilateral triangle's area, one would get the final result of $\sqrt{300} + \sqrt{252}$.

 This makes the answer $252 + 300 = 552$. $\boxed{\textbf{D}.}$




\textbf{Solution 4}


                    \begin{figure}[H]
                        \centering
                        \includegraphics[width=0.5\textwidth]{Solution/AMC_12_12B_2018_25_Solution_Figure_1.png}
                    \end{figure}
                    

 First, note that because the $\angle P_1=\angle P_2=\angle P_3=\pi/3$, the arcs inside the shaded equilateral triangle are each $2\pi/3$. Also, the distances between the centers of any two of the $3$ given circles are each $8$.Draw the circle $\Gamma$ concentric with $\omega_1$ with radius $2$. Because the arc of $\omega_1$ inside said triangle is $2\pi/3$, $\Gamma$ touches $P_1P_3$, say at a point $X$. Thus, $P_1P_3$ is a common tangent of $\omega_3$ and $\Gamma$, and it can be seen from inspection of the given diagram that the line is an common internal tangent.The length of the common internal tangent segment $XP_3$ of $\Gamma$ and $\omega_3$ is then $\sqrt{8^2-(2+4)^2}=2\sqrt{7}$, and it is easily seen that $XP_1=4\sin \pi/3=2\sqrt{3}$.Because $P_1P_3=2(\sqrt{3}+\sqrt{7})$, the area of the shaded equilateral triangle is $\sqrt{3}(\sqrt{3}+\sqrt{7})^2=10\sqrt{3}+6\sqrt{7}$. We get $\sqrt{300}+\sqrt{252}\Rightarrow\boxed{\textbf{(D) }552}.$




\textbf{Solution 5 (Pure Coordinate Bash)}

This seems like a coordinate bashable problem. To do this, we notice that it is easier to graph the equilateral triangle first, then the circles, rather than the other way around. Let's forget the lengths of the radius for a moment and focus instead on the ratio of the circles' radius to $\triangle P_1P_2P_3$'s side length.

 WLOG let $P_1P_2 = 2$. Place triangle $P_1P_2P_3$ on the coordinate plane with $P_1(0,\sqrt3)$, $P_2(1,0)$, and $P_3(-1,0)$. Let $r$ be the radius of the circles. 

 Now, we find the coordinates of the centers of circles $\omega_1$ and $\omega_2$. Since $\omega_2$ is tangent to the x-axis at $(1,0)$, the center of $\omega_2$ is $(1,r)$. Draw a right triangle with legs parallel to the x and y axes, and with hypotenuse as the segment from the center of $\omega_1$ to $P_1$. Since the slope of $\overline{P_1P_2}$ is $-\sqrt3$, the slope of the hypotenuse is $\frac{1}{\sqrt{3}}$, so the right triangle is $30-60-90$. It's easy to see that the center of $\omega_1$ is $\left(-\frac{r\sqrt{3}}{2}, \sqrt{3} - \frac{r}{2}\right)$.

 Since $\omega_1$ and $\omega_2$ are tangent, the distance between the centers is $2r$, so we have
 \[\sqrt{\left(1 + \frac{r\sqrt{3}}{2}\right)^2 + \left(\frac{3r}{2} - \sqrt{3}\right)^2} = 2r\]
 \[\left(1 + \frac{r\sqrt{3}}{2}\right)^2 + \left(\frac{3r}{2} - \sqrt{3}\right)^2 = 4r^2\]
 \[1 + r\sqrt{3} + \frac{3r^2}{4} + \frac{9r^2}{4} - 3r\sqrt{3} + 3 = 4r^2\]
 \[4 - 2r\sqrt{3} + 3r^2 = 4r^2\]
 \[r^2 + 2r\sqrt{3} - 4 = 0\]By the Quadratic Formula, $r = \frac{-2\sqrt{3} \pm \sqrt{12 + 16}}{2} = -\sqrt{3}\pm\sqrt{7}$.We take the positive value to get the radius of the circle is $\sqrt{7} - \sqrt{3}$.

 Therefore, the ratio of the radius to the side length of the equilateral triangle is $\sqrt{7} - \sqrt{3} : 2 = 4 : \frac{8}{\sqrt{7} - \sqrt{3}} = 4 : 2(\sqrt{7} + \sqrt{3})$. 

 The side length of the equilateral triangle is $2(\sqrt{7} + \sqrt{3})$, so its area is $(\sqrt{7} + \sqrt{3})^2\sqrt{3} = 10\sqrt{3} + 6\sqrt{7} = \sqrt{300} + \sqrt{252} \implies 300 + 252 = \boxed{\textbf{(D) }552}$.

 



\textbf{Solution 6 (Trignometry Bash)}

Let $O_1$, $O_2$, and $O_3$ be the centers of $\omega_1$, $\omega_2$, and $\omega_3$

 Let $P_1P_2 = a$, $\angle O_2O_1P_1= \alpha$

 \[[O_1P_1O_2P_2O_3P_3] = [P_1P_2P_3] + 3 \cdot [O_1P_1P_3] = [O_1O_2O_3] + 3 \cdot [O_1P_1O_2]\]

 \[[P_1P_2P_3] + 3 \cdot [O_1P_1P_3] = \frac12 \cdot a \cdot \frac{ \sqrt{3} }{2} a + 3 \cdot \frac12 \cdot 4 \cdot a \cdot \sin 30^{\circ} = \frac{ \sqrt{3} }{4} a^2 + 3a\]

 \[[O_1O_2O_3] + 3 \cdot [O_1P_1O_2] = \frac12 \cdot 8 \cdot \frac{ \sqrt{3} }{2} \cdot 8 + 3 \cdot \frac12 \cdot 4 \cdot 8 \cdot \sin \alpha = 16 \sqrt{3} + 48 \sin \alpha\]

 \[\frac{ \sqrt{3} }{4} a^2 + 3a = 16 \sqrt{3} + 48 \sin \alpha\]

 \[\sin \alpha = \frac{ \sqrt{3} }{192} a^2 + \frac{a}{16} - \frac{ \sqrt{3} }{3}\]

 By the Law of Cosine from $\triangle O_2P_1P_2$, $O_2P_1^2 = a^2 + 4^2 - 2 \cdot a \cdot 4 \cdot \cos 30^{\circ} = a^2 + 16 - 4a \sqrt{3}$

 By the Law of Cosine from $\triangle O_1O_2P_1$, $O_2P_1^2 = 4^2 + 8 ^ 2 - 2 \cdot 4 \cdot 8 \cdot \cos \alpha = 80 - 64 \cos \alpha$

 \[a^2 + 16 - 4a \sqrt{3} = 80 - 64 \cos \alpha\]

 \[\cos \alpha = -(\frac{a^2}{64} - \frac{ \sqrt{3} }{16} a - 1)\]

 \[\because \sin^2 \alpha + \cos^2 \alpha = 1\]

 \[\therefore (\frac{ \sqrt{3} }{192} a^2 + \frac{a}{16} - \frac{ \sqrt{3} }{3})^2 + (\frac{a^2}{64} - \frac{ \sqrt{3} }{16} a - 1)^2 = 1\]

 \[[ \frac{1}{\sqrt{3}} (\frac{ a^2 }{64}  + \frac{ \sqrt{3} }{16} a - 1)]^2 + (\frac{a^2}{64} - \frac{ \sqrt{3} }{16} a - 1)^2 = 1\]

 \[(\frac{ a^2 }{64} + \frac{ \sqrt{3} }{16} a - 1)^2 + 3 \cdot (\frac{a^2}{64} - \frac{ \sqrt{3} }{16} a - 1)^2 = 3\]

 \[(a^2 + 4 \sqrt{3} a - 64)^2 + 3 \cdot (a^2 - 4 \sqrt{3} a - 64)^2 = 3 \cdot 64^2\]

 \[(a^2 - 64)^2 + 2 (a^2 - 64) \cdot 4 \sqrt{3} a + (4 \sqrt{3} a)^2 + 3(a^2 - 64)^2 - 6 (a^2 - 64) \cdot 4 \sqrt{3} a + 3(4 \sqrt{3} a)^2 = 3 \cdot 64^2\]

 \[4 (a^2 - 64)^2 - 4 (a^2 - 64) \cdot 4 \sqrt{3} a + 192 a^2 - 3 \cdot 64^2 = 0\]

 \[4 (a^2 - 64)(a^2 - 64 - 4 \sqrt{3} a) + 192(a^2 - 64) = 0\]

 \[4 (a^2 - 64)(a^2 - 64 - 4 \sqrt{3} a + 48) = 0\]

 \[(a^2 - 64)(a^2 - 4 \sqrt{3} a - 16) = 0\]

 \[a^2 -64 = 0 \quad or \quad a^2 - 4 \sqrt{3} a - 16 = 0\]

 If $a^2 =64$, $a= 8$, $P_1O_2^2 = 8^2 + 4^2 - 2 \cdot 8 \cdot 4 \cdot \cos 30^{\circ} = 80 - 32 \sqrt{3} < 16$, $P_1O_2 < 4$, meaning that $P_1$ is inside the circle.

 However, $P_1$ is not inside the circle.

 
 \[So, \quad a^2 - 4 \sqrt{3} a - 16 = 0\]

 \[a = \frac{ 4 \sqrt{3} + \sqrt{ (4 \sqrt{3})^2 - 4(-16) } }{2} = 2 \sqrt{3} + 2 \sqrt{7}\]

 \[[P_1P_2P_3] = a^2 \cdot \frac{ \sqrt{3} }{4} = (2 \sqrt{3} + 2 \sqrt{7})^2 \cdot \frac{ \sqrt{3} }{4} = \sqrt{300} + \sqrt{252}\]

 Therefore, the answer is $300 + 252 = \boxed{\textbf{(D) } 552}$.

 



\textbf{Solution 7}


                    \begin{figure}[H]
                        \centering
                        \includegraphics[width=0.5\textwidth]{Solution/AMC_12_12B_2018_25_Solution_Figure_2.png}
                    \end{figure}
                    

 Suppose $P_1P_2 = x$, noticed that $OMP_1H$ is rectangle.
 \[OO_1 = \frac{O_1O_2}{\sqrt{3}} = \frac83\sqrt{3}\]
 \[OH = P_1M = \frac{1}{2}P_1P_2 = \frac12 x\]
 \[O_1H = O_1P_1 - HP_1 = O_1P_1 - OM = O_1P_1 - \frac{P_1P_2}{2\sqrt{3}} = 4 - \frac{\sqrt{3}}{6}x\]$\triangle O_1HO$ is right triangle, we can use Pythagorean theorem to establish an equation,
 \[OH^2 + O_1H^2 = OO_1^2\]
 \[\left(\frac12 x\right)^2 + \left(4 - \frac{\sqrt{3}}{6}x\right)^2 = \left(\frac83\sqrt{3}\right)^2\]
 \[x^2 - 4\sqrt{3}x - 16 = 0\]
 \[x = 2\sqrt{3} + 2\sqrt{7}\]The side length of the equilateral triangle is $2(\sqrt{7} + \sqrt{3})$, so its area is $(\sqrt{7} + \sqrt{3})^2\sqrt{3} = 10\sqrt{3} + 6\sqrt{7} = \sqrt{300} + \sqrt{252} \implies 300 + 252 = \boxed{\textbf{(D) }552}$.

 



\textbf{Solution 8 (Vector Method)}

Let $x$ be the side length of the equilateral $\triangle P_1P_2P_3$. We analyze the vector path from $O_1$ to $O_2$:
 \[\vec{O_1O_2} = \vec{O_1P_1} + \vec{P_1P_2} + \vec{P_2O_2}\]

 We know the magnitudes of these vectors:

 
            \begin{itemize}
        \item $|\vec{O_1O_2}| = 4 + 4 = 8$ (since the circles are externally tangent)
        \item $|\vec{O_1P_1}| = 4$ (radius of $\omega_1$)
        \item $|\vec{P_1P_2}| = x$
        \item $|\vec{P_2O_2}| = 4$ (radius of $\omega_2$)
            \end{itemize}
            Next, we determine the pairwise angles between these vectors to evaluate their dot products:

 1. \textbf{Angle between $\vec{O_1P_1}$ and $\vec{P_1P_2}$}:Since $P_1P_2$ is tangent to $\omega_1$ at $P_1$, the radius is perpendicular to the tangent line. Thus, the angle is $90^\circ$, and their dot product is:
 \[\vec{O_1P_1} \cdot \vec{P_1P_2} = 0\]

 2. \textbf{Angle between $\vec{P_1P_2}$ and $\vec{P_2O_2}$}:In equilateral $\triangle P_1P_2P_3$, the interior angle $\angle P_1P_2P_3 = 60^\circ$. Since $P_2P_3$ is tangent to $\omega_2$ at $P_2$, we have $O_2P_2 \perp P_2P_3$. This implies that $\angle O_2P_2P_1 = 90^\circ - 60^\circ = 30^\circ$.The vector $\vec{P_2P_1}$ forms a $30^\circ$ angle with $\vec{P_2O_2}$. Since $\vec{P_1P_2}$ points in the exact opposite direction of $\vec{P_2P_1}$, the angle between $\vec{P_1P_2}$ and $\vec{P_2O_2}$ is $180^\circ - 30^\circ = 150^\circ$. Their dot product is:
 \[\vec{P_1P_2} \cdot \vec{P_2O_2} = x \cdot 4 \cdot \cos(150^\circ) = -2\sqrt{3}x\]

 3. \textbf{Angle between $\vec{O_1P_1}$ and $\vec{P_2O_2}$}:The vector $\vec{O_1P_1}$ is the outward-pointing normal to the side $P_1P_2$, while $\vec{O_2P_2}$ is the outward-pointing normal to the side $P_2P_3$. The angle between the outward normals of two adjacent sides of an equilateral triangle is equal to its exterior angle, which is $120^\circ$. Since $\vec{P_2O_2} = -\vec{O_2P_2}$ (it is the inward-pointing normal relative to the triangle's corner), the angle between $\vec{O_1P_1}$ and $\vec{P_2O_2}$ is $180^\circ - 120^\circ = 60^\circ$. Their dot product is:
 \[\vec{O_1P_1} \cdot \vec{P_2O_2} = 4 \cdot 4 \cdot \cos(60^\circ) = 16 \cdot \frac{1}{2} = 8\]

 Now, we square both sides of the vector equation:
 \[|\vec{O_1O_2}|^2 = |\vec{O_1P_1}|^2 + |\vec{P_1P_2}|^2 + |\vec{P_2O_2}|^2 + 2(\vec{O_1P_1} \cdot \vec{P_1P_2}) + 2(\vec{P_1P_2} \cdot \vec{P_2O_2}) + 2(\vec{O_1P_1} \cdot \vec{P_2O_2})\]

 Substitute the known magnitudes and dot products into the equation:
 \[8^2 = 4^2 + x^2 + 4^2 + 2(0) + 2(-2\sqrt{3}x) + 2(8)\]
 \[64 = 16 + x^2 + 16 - 4\sqrt{3}x + 16\]
 \[64 = x^2 - 4\sqrt{3}x + 48\]

 Rearranging this into a standard quadratic equation yields:
 \[x^2 - 4\sqrt{3}x - 16 = 0\]

 Using the quadratic formula to solve for $x$ (and taking the positive root since $x$ represents a length):
 \[x = \frac{4\sqrt{3} + \sqrt{(-4\sqrt{3})^2 - 4(1)(-16)}}{2} = \frac{4\sqrt{3} + \sqrt{48 + 64}}{2} = \frac{4\sqrt{3} + \sqrt{112}}{2}\]
 \[x = \frac{4\sqrt{3} + 4\sqrt{7}}{2} = 2\sqrt{3} + 2\sqrt{7}\]

 Finally, we calculate the area of the equilateral $\triangle P_1P_2P_3$:
 \[[P_1P_2P_3] = \frac{\sqrt{3}}{4} x^2 = \frac{\sqrt{3}}{4} (2\sqrt{3} + 2\sqrt{7})^2 = \sqrt{300} + \sqrt{252}\]

 Therefore, the requested sum is $300 + 252 = \boxed{\textbf{(D) } 552}$.

 


\end{enumerate}